\documentclass[letterpaper, 10 pt, conference]{ieeeconf}
% NOTE: This template will need to change to the target journal's LaTeX class
% before submission (e.g., IEEEtran for IEEE Access, elsarticle for Elsevier,
% svjour3 for Springer). The content below is journal-ready.
%
% -----------------------------------------------------------------------------
% FOUR-DATASET EXPANSION -- STATUS:
%   - Datasets subsection: lists all 4 (BC, HD, Pima, ILPD)          [DONE]
%   - Methods/run-count updated (360 -> 720 runs)                    [DONE]
%   - Pima and ILPD citations in bibliography                        [DONE]
%   - Tables tab:pima_results / tab:ilpd_results                     [DONE]
%   - Cross-dataset summary table rows for Pima/ILPD                 [DONE]
%   - Pima and ILPD interpretive prose (RF / SVM / LR per dataset)   [DONE]
%   - Abstract rewritten for 4-dataset scope with ILPD/SVM narrative [DONE]
%   - Cross-dataset synthesis paragraph (4 patterns)                 [DONE]
%   - Contribution count updated (17/18 -> 3/72 cells)               [DONE]
%
%   - Conclusion rewritten with 6-finding enumeration + deployment advice [DONE]
%   - Tables refreshed with 10-seed 4-dataset numbers (BC/HD/PD/ILPD)    [DONE]
%   - Figures regenerated from results_raw.csv (4 rows x 3 cols facet)  [DONE]
%   - Interpretive prose refreshed; HD/RF p=0.016 claim retracted       [DONE]
%   - Summary table + synthesis paragraph updated (Holm-corrected 0/60) [DONE]
%
% Remaining before submission:
%   1. Template swap to IEEEtran (or elsarticle / svjour3) for target journal.
%   2. Populate real GitHub repo URL in Reproducibility section.
%   3. Run supplementary Colab for (a)/(b) panels of fig:meta_analysis.
%   4. Supervisor review; cover letter.
% -----------------------------------------------------------------------------

\IEEEoverridecommandlockouts
\overrideIEEEmargins

% Packages
\usepackage{graphicx}
\graphicspath{{./}{./figures/}}
\usepackage{amsmath}
\usepackage{amssymb}
\usepackage{booktabs}
\usepackage{multirow}
\usepackage{array}

\title{\LARGE \bf
An Experimental Comparison of Hyperparameter Optimization Methods for Healthcare Machine Learning: From Grid Search to Swarm Intelligence
}

\author{Saadeddin Hroub$^{1}$, Masri Ayob$^{2}$, and Mohd Nor Akmal Khalid$^{3}$%
	\thanks{*This work was supported by the Ministry of Higher Education (MoHE) Malaysia under the Transdisciplinary Research Grant Scheme (FRGS/1/2024/ICT02/UKM/01/).}%
	\thanks{$^{1,2,3}$The authors are with the Data Mining and Optimization Lab, Faculty of Information Science and Technology, Universiti Kebangsaan Malaysia, Selangor, Malaysia.}%
	\thanks{$^{1}${\tt\small saad.hroub@gmail.com}}%
	\thanks{$^{2}${\tt\small masri@ukm.edu.my} (Corresponding author)}%
	\thanks{$^{3}${\tt\small akmal@ukm.edu.my}}%
}

\begin{document}
	\raggedbottom
	\maketitle
	\thispagestyle{empty}
	\pagestyle{empty}

	%%%%%%%%%%%%%%%%%%%%%%%%%%%%%%%%%%%%%%%%%%%%%%%%%%%%%%%%%%%%%%%%%%%%%%%%%%%%%%%%
	\begin{abstract}

The performance of machine learning models in healthcare depends strongly on hyperparameter optimization (HPO), yet current practice relies disproportionately on grid search or default settings, which are slow, inconsistent, or unsuitable for clinical environments that demand fast, reliable, and resource-efficient optimization. Existing comparative studies are fragmented: healthcare-oriented work typically employs traditional methods without meta-heuristic baselines, while meta-heuristic studies rarely benchmark against traditional methods or manual tuning. This paper addresses the gap with a systematic, head-to-head evaluation of six HPO strategies---manual tuning, grid search, random search, Bayesian optimization, Grey Wolf Optimizer (GWO), and Particle Swarm Optimization (PSO)---across three classifiers (Random Forest, Support Vector Machine, Logistic Regression) on four benchmark healthcare datasets spanning oncology, cardiology, endocrinology, and hepatology (Breast Cancer Wisconsin, Heart Disease Statlog, Pima Indians Diabetes, Indian Liver Patient Dataset), using ten-run statistical validation with the Wilcoxon signed-rank test and Holm--Bonferroni multiple-comparison correction. Across the $60$ non-manual model--method--dataset comparisons in the benchmark, meta-heuristics attained the numerically best accuracy in only one and produced \emph{zero} Holm-corrected significant accuracy improvements over the simplest-available alternative; the only Holm-corrected significant departures from manual tuning are two accuracy \emph{decreases} on ILPD with an SVM classifier. Random search emerged as the Pareto-dominant optimizer, reaching the same accuracy plateau as GWO and PSO at $3$--$25\times$ lower cost on every dataset; Bayesian optimization matched random search's accuracy at a comparable budget. The single cell where meta-heuristics earn their additional compute is the Indian Liver Patient Dataset with an SVM classifier, where GWO and PSO escape a degenerate majority-class basin and improve macro-$F_1$ from $0.4164$ (manual baseline) to $\approx 0.565$ at a $2.7$-point accuracy cost---a trade-off that is favourable for clinical deployment but invisible under accuracy-only evaluation. We attribute the systematic absence of meta-heuristic superiority to strong performance plateaus, mixed discrete--continuous search spaces, and validation-set noise on small clinical datasets, and we provide evidence-based deployment guidelines that link dataset size, model complexity, and class imbalance to the most appropriate HPO strategy for clinical use.

	\end{abstract}

	\textbf{Index Terms}---Hyperparameter optimization, Random Forest, swarm intelligence, healthcare machine learning, meta-heuristic algorithms, clinical deployment.

	%%%%%%%%%%%%%%%%%%%%%%%%%%%%%%%%%%%%%%%%%%%%%%%%%%%%%%%%%%%%%%%%%%%%%%%%%%%%%%%%
	\section{INTRODUCTION}
	\label{sec:intro}

\subsection{Motivation and Background}

The predictive performance of machine learning models depends critically on their hyperparameter configuration---settings that control the learning process but cannot be learned from data directly~\cite{probst2019}. For Random Forest (RF) models, hyperparameters such as the number of trees, maximum depth, minimum samples per split, and feature subset size directly affect decision boundaries, generalization capability, and computational efficiency~\cite{probst2019, breiman2001}. The choice of HPO strategy is therefore a fundamental step in machine-learning deployment, and nowhere is that more important than in high-stakes healthcare domains, where suboptimal configurations can directly impact patient outcomes~\cite{rahimi2016}.

Healthcare applications impose constraints that generic HPO approaches systematically fail to meet. Clinical workflows demand model deployment within hours, not the one-to-three days required by exhaustive optimization~\cite{elgeldawi2021, bischl2023}. Patient safety requires consistent, reproducible performance across diverse clinical settings, which rules out methods that are highly sensitive to initialization or exhibit large run-to-run variance. Healthcare stakeholders require transparent, explainable optimization processes; manual tuning lacks this transparency, and black-box Bayesian methods obscure it~\cite{naik2022}. Finally, evolving patient demographics, treatment protocols, and resource constraints require continuous re-optimization that static HPO frameworks cannot provide~\cite{elgeldawi2021}.

\subsection{Limitations of Existing Methods}

Despite extensive research on HPO, current approaches exhibit systematic limitations when applied to healthcare classification. \emph{Manual tuning} is computationally cheap but non-reproducible; different practitioners derive incompatible configurations for the same clinical problem, undermining reproducibility~\cite{bischl2023, raji2022}. \emph{Grid search} suffers from exponential cost $\mathcal{O}(n^d)$ and becomes prohibitive once $6$--$8$ hyperparameters are in play~\cite{elgeldawi2021, bischl2023}. \emph{Random search} is cheaper than grid search but carries no learning mechanism across evaluations, wasting budget on poor configurations~\cite{bergstra2012}. \emph{Bayesian optimization} is sample-efficient but relies on Gaussian-process assumptions that are often inadequate for mixed discrete--continuous spaces typical of healthcare HPO~\cite{abbas2023}.

Meta-heuristic algorithms---Particle Swarm Optimization (PSO), Genetic Algorithms (GA), Grey Wolf Optimizer (GWO), and Whale Optimization Algorithm (WOA)---provide gradient-free search suitable for Random Forest hyperparameter spaces~\cite{huo2022, otero2012, tayebi2022}. However, standard implementations suffer from premature convergence, with populations losing diversity in early iterations~\cite{huo2022, yang2023}. This diversity collapse prevents escape from local optima in the complex, multimodal landscapes characteristic of healthcare datasets. Recent studies also report a striking \emph{performance plateau} effect: in healthcare ML, optimized configurations often provide only marginal improvements over default settings, suggesting that conventional HPO methods stagnate in local optima without effective mechanisms to escape them.

\subsection{Research Question and Contributions}

These observations motivate the central question of the present study: \emph{do meta-heuristic HPO methods outperform traditional methods on healthcare classification tasks, and if not, why not?} To answer this, we conduct a unified, head-to-head comparison of six HPO methods spanning manual, traditional, Bayesian, and meta-heuristic families on four benchmark healthcare datasets spanning oncology, cardiology, endocrinology, and hepatology, using ten-run statistical validation with the Wilcoxon signed-rank test and Holm--Bonferroni multiple-comparison correction.

The main contributions of this work are as follows.
\begin{enumerate}
	\item \textbf{Unified empirical comparison.} We provide the first systematic evaluation---to our knowledge---that includes manual tuning, grid search, random search, Bayesian optimization, GWO, and PSO within a single experimental framework on healthcare classification. Each configuration is evaluated across ten independent runs with formal statistical significance testing.
	\item \textbf{Evidence against meta-heuristic superiority in this regime.} Across the $60$ non-manual comparisons, GWO and PSO attain the numerically best accuracy in only one cell (BC / RF / GWO, $p = 0.516$) and produce zero Holm-corrected significant improvements over the simplest alternative; the only Holm-corrected significant departures from manual tuning are two accuracy \emph{decreases}. These findings contradict claims common in the recent HPO literature. We provide mechanistic explanations grounded in landscape characteristics, population dynamics, and validation-set noise.
	\item \textbf{Healthcare-specific constraint analysis.} We formalize temporal, interpretability, reproducibility, fairness, computational, and adaptation constraints that distinguish healthcare HPO from generic benchmarks, and we quantify how each method fares against them.
	\item \textbf{Deployment guidelines.} We derive evidence-based guidelines linking dataset size and model complexity to the most appropriate HPO strategy for clinical use, together with a decision framework for method selection under time and resource constraints.
	\item \textbf{Research agenda.} We outline a concrete research programme with testable hypotheses for enhanced swarm intelligence, hybrid Bayesian--swarm optimization, noise-aware fitness, and transfer learning across institutions.
\end{enumerate}

\subsection{Paper Organization}

The remainder of the paper is organized as follows. Section~\ref{sec:related} reviews related work and identifies the research gap. Section~\ref{sec:hpo} presents the HPO methods evaluated. Section~\ref{sec:healthcare} analyses healthcare-specific requirements. Section~\ref{sec:methodology} describes the experimental methodology. Section~\ref{sec:results} reports empirical results. Section~\ref{sec:discussion} discusses why meta-heuristics underperformed and provides clinical deployment implications. Section~\ref{sec:limitations} acknowledges study limitations. Section~\ref{sec:future} outlines future research directions with testable hypotheses. Section~\ref{sec:conclusion} concludes.

%%%%%%%%%%%%%%%%%%%%%%%%%%%%%%%%%%%%%%%%%%%%%%%%%%%%%%%%%%%%%%%%%%%%%%%%%%%%%%%%
\section{RELATED WORK AND RESEARCH GAP}
\label{sec:related}

The HPO literature has grown rapidly but splits into two largely disjoint streams: studies that apply traditional methods in domain contexts (including healthcare), and studies that propose or refine meta-heuristic methods. A comprehensive comparison that bridges the two streams, includes a manual baseline, and reports statistical validation on healthcare tasks has been notably absent.

\subsection{Traditional HPO Methods in Healthcare}

Several healthcare-oriented studies have applied traditional HPO methods to classification tasks. Gao and Ding~\cite{gao2020} compared grid search, random search, and Bayesian optimization for disease prediction on breast cancer and cardiovascular datasets, finding Bayesian optimization more stable but excluding meta-heuristics entirely. Ebrahim et al.~\cite{ebrahim2024} employed grid search for breast cancer recurrence prediction across multiple classifiers including logistic regression, decision trees, and deep neural networks, but did not systematically compare alternative HPO strategies. Zhou et al.~\cite{zhou2024} used grid search for an AdaBoost--logistic breast cancer classifier, achieving $99.12\%$ accuracy on their dataset but without exploring whether alternative optimization methods could improve results. Hassan et al.~\cite{hassan2023} focused on LASSO-based feature selection for breast cancer detection, comparing classifier performance but not HPO methodologies. Jain et al.~\cite{jain2024} performed parametric optimization for breast cancer diagnosis comparing ML and deep-learning algorithms, yet did not systematically evaluate different HPO strategies. Across these studies, traditional methods---predominantly grid search---are employed without considering whether meta-heuristic alternatives might offer advantages.

\subsection{Meta-heuristic HPO Approaches}

Conversely, meta-heuristic HPO studies typically exclude traditional baselines. Tayebi and El Kafhali~\cite{tayebi2022} provided a broad meta-heuristic comparison, evaluating GA, DE, ABC, GWO, PSO, and TLBO against grid search for fraud detection; their work excluded Bayesian optimization, random search, and a manual baseline, and targeted fraud detection rather than healthcare. Ulutas et al.~\cite{ulutas2024} developed a PSO--GWO hybrid for diabetes detection using ensemble learning, comparing only meta-heuristic variants (PSO, GWO, PSO--GWO) without traditional baselines. Kilicarslan~\cite{kilicarslan2023} applied a PSO+GWO hybrid for CNN-based cardiovascular disease detection but focused exclusively on deep learning architectures without traditional HPO comparison. Challapalli and Devarakonda~\cite{challapalli2022} proposed a hybrid PSO--GWO for CNN image classification in a non-healthcare domain without traditional baselines. Mohakud and Dash~\cite{mohakud2022} used an enhanced GWO for skin-cancer image segmentation with deep learning, comparing only meta-heuristic variants. Alibrahim and Ludwig~\cite{alibrahim2021} compared GA against grid search and Bayesian optimization for neural-network tuning, providing useful insight but excluding PSO, GWO, and healthcare applications. Vincent and Jidesh~\cite{vincent2023} surveyed evolutionary algorithms for AutoML hyperparameter optimization, offering a methodological framework without domain-specific healthcare validation.

\subsection{Research Gap}

Table~\ref{tab:related_work} synthesizes the literature. No existing study provides a unified comparison of traditional methods (grid search, random search, Bayesian optimization) alongside meta-heuristics (GWO, PSO) with a manual baseline on healthcare classification tasks using rigorous statistical validation. The fragmentation creates four concrete problems: (i)~practitioners cannot determine whether sophisticated meta-heuristics justify their computational overhead compared to simpler random search; (ii)~the absence of a manual baseline obscures whether any optimization provides genuine benefit over expert defaults; (iii)~the lack of statistical validation across multiple runs undermines the reliability of single-run comparisons; and (iv)~healthcare-specific metrics such as the false negative rate (FNR)---critical in clinical settings where missed diagnoses carry severe consequences---are almost never reported. The present study is designed to fill each of these gaps.

\begin{table*}[t]
	\caption{Coverage of HPO methods, statistical validation, and healthcare applicability in representative related studies.}
	\label{tab:related_work}
	\centering
	\scriptsize
	\begin{tabular}{|p{2.5cm}|c|c|c|c|c|c|c|c|}
		\hline
		\textbf{Study} & \textbf{Grid} & \textbf{Random} & \textbf{Bayes.} & \textbf{GWO} & \textbf{PSO} & \textbf{Manual} & \textbf{Stats.} & \textbf{Health.} \\
		\hline
		Tayebi~\cite{tayebi2022} & \checkmark & -- & -- & \checkmark & \checkmark & -- & -- & -- \\
		Gao~\cite{gao2020} & \checkmark & \checkmark & \checkmark & -- & -- & -- & -- & \checkmark \\
		Ulutas~\cite{ulutas2024} & -- & -- & -- & \checkmark & \checkmark & -- & -- & \checkmark \\
		Ebrahim~\cite{ebrahim2024} & \checkmark & -- & -- & -- & -- & -- & -- & \checkmark \\
		Kilicarslan~\cite{kilicarslan2023} & -- & -- & -- & \checkmark & \checkmark & -- & -- & \checkmark \\
		Alibrahim~\cite{alibrahim2021} & \checkmark & -- & \checkmark & -- & -- & -- & -- & -- \\
		Challapalli~\cite{challapalli2022} & -- & -- & -- & \checkmark & \checkmark & -- & -- & -- \\
		Mohakud~\cite{mohakud2022} & -- & -- & -- & \checkmark & \checkmark & -- & -- & \checkmark \\
		\hline
		\textbf{This study} & \checkmark & \checkmark & \checkmark & \checkmark & \checkmark & \checkmark & \checkmark & \checkmark \\
		\hline
	\end{tabular}
\end{table*}

%%%%%%%%%%%%%%%%%%%%%%%%%%%%%%%%%%%%%%%%%%%%%%%%%%%%%%%%%%%%%%%%%%%%%%%%%%%%%%%%
\section{HYPERPARAMETER OPTIMIZATION METHODS}
\label{sec:hpo}

This section describes the six HPO methods evaluated in this study, together with their advantages, disadvantages, and implications for healthcare deployment.

\subsection{Manual Tuning}

Manual tuning relies on domain expertise, trial-and-error experimentation, and iterative refinement~\cite{bischl2023}. Practitioners leverage understanding of algorithm behaviour, dataset characteristics, and task requirements to select and adjust hyperparameters. Its strengths are computational efficiency, the incorporation of domain knowledge, rapid prototyping, and interpretability of decisions. Its limitations are non-reproducibility across practitioners, time-intensive iteration on complex models, subjectivity, inconsistency, failure to systematically explore large search spaces, and vulnerability to cognitive biases and local optima. In healthcare, where reproducibility and regulatory compliance are mandatory, different practitioners can derive incompatible configurations for the same clinical problem, undermining model standardization.

\subsection{Grid Search}

Grid search exhaustively evaluates all combinations within a predefined hyperparameter grid~\cite{elgeldawi2021}. For each hyperparameter, practitioners specify discrete candidate values, and the method trains and evaluates a model for every combination, selecting the configuration that yields the best validation performance. Grid search is reproducible, simple to implement, and easily parallelizable, and it guarantees optimality within the chosen grid. However, its computational cost grows exponentially with dimensionality, $\mathcal{O}(n^d)$, and it handles continuous hyperparameters inefficiently. In practice, grid search becomes prohibitive for Random Forest optimization requiring simultaneous tuning of $6$--$8$ hyperparameters, and clinical deployment timelines are typically incompatible with exhaustive search.

\subsection{Random Search}

Random search samples hyperparameter configurations from specified distributions and evaluates a fixed number of random combinations~\cite{bergstra2012}. Bergstra and Bengio~\cite{bergstra2012} showed that random search is more efficient than grid search in high-dimensional optimization where only a few hyperparameters dominate performance, because random sampling provides better coverage of important dimensions at equivalent budget. Random search is computationally efficient, scales well with dimensionality, is easily parallelizable, and handles continuous and discrete parameters naturally. Its main drawback is the absence of a learning mechanism: evaluations do not inform subsequent sampling, so the method cannot exploit promising regions identified during the search. Healthcare deployment benefits from random search's simplicity and reproducibility (given fixed seeds), though its lack of learning is sometimes cited as a disadvantage.

\subsection{Bayesian Optimization}

Bayesian optimization (BO) treats HPO as a sequential decision problem, constructing a probabilistic surrogate model to guide the search toward promising configurations~\cite{snoek2012, feurer2019}. The method alternates between updating a probabilistic model of the objective function and selecting the next configuration to evaluate based on an acquisition function.

The surrogate is typically a Gaussian process (GP) modelling the relationship between hyperparameters and validation performance; the GP provides predictions with uncertainty estimates for unexplored configurations. The acquisition function guides the exploration--exploitation trade-off; common choices include Expected Improvement, Probability of Improvement, and Upper Confidence Bound. In each iteration, the method selects the configuration that maximizes the acquisition function, evaluates actual performance, and updates the surrogate with the new observation.

BO is sample-efficient, incorporates uncertainty, balances exploration and exploitation, and handles expensive evaluations effectively. Its limitations include the inadequacy of GP assumptions for complex non-linear relationships, poor scaling to high-dimensional spaces (beyond roughly $20$ parameters), overhead from surrogate-model updates, sensitivity to kernel choice, and difficulty handling discrete--continuous mixed spaces common in Random Forest tuning. In healthcare contexts, the black-box nature of BO conflicts with interpretability requirements, and its limited theoretical guarantees can be problematic for safety-critical deployment. In this work, we use the Tree-structured Parzen Estimator (TPE) implementation via Optuna, which handles mixed spaces more gracefully than standard GP-based BO.

\subsection{Meta-heuristic Algorithms}

Meta-heuristics provide gradient-free optimization via nature-inspired search~\cite{alyasseri2020}, maintaining a population of candidate solutions that is refined iteratively through mechanisms inspired by biological, physical, or social phenomena.

\subsubsection{Particle Swarm Optimization (PSO)}

PSO simulates the social behaviour of bird flocking or fish schooling~\cite{rehan2022, kennedy1995}. Each particle represents a candidate hyperparameter configuration and updates its velocity and position based on its personal best and the swarm-wide best:
\begin{align}
	v_i(t+1) &= w\,v_i(t) + c_1 r_1 (p_i - x_i(t)) + c_2 r_2 (g - x_i(t)), \\
	x_i(t+1) &= x_i(t) + v_i(t+1),
\end{align}
where $v_i$ and $x_i$ denote the velocity and position of particle $i$, $p_i$ is its personal best position, $g$ is the global best position, $w$ is the inertia weight, $c_1, c_2$ are acceleration coefficients, and $r_1, r_2 \in [0,1]$ are uniform random scalars drawn at each step. PSO is simple to implement, has few control parameters, and exhibits strong global-search capability, but it is well known to suffer from premature convergence once particles cluster around a local optimum and from velocity explosion in unbounded spaces.

\subsubsection{Grey Wolf Optimizer (GWO)}

GWO models the hunting hierarchy of grey wolves~\cite{huo2022, mirjalili2014}. Wolves maintain a social hierarchy---alpha, beta, delta, and omega---with positions updated based on the three best solutions. The update rule can be written
\begin{align}
	D_{\alpha} &= |C_1 X_{\alpha} - X|, \\
	X(t+1) &= (X_1 + X_2 + X_3)/3,
\end{align}
where $X_{\alpha}, X_{\beta}, X_{\delta}$ denote the top three solutions, $C_1$ is a coefficient vector that controls the leader's influence, and $X_1, X_2, X_3$ are intermediate candidate positions derived analogously from $X_{\alpha}, X_{\beta}, X_{\delta}$ respectively (the symmetric expressions for $D_{\beta}, D_{\delta}$ and for $X_2, X_3$ follow the same pattern and are omitted for brevity; see~\cite{mirjalili2014}). GWO has few control parameters and converges quickly, but its exploration ability deteriorates in later iterations as population diversity collapses and wolves concentrate around the alpha position.

\subsubsection{Other Nature-inspired Meta-heuristics}

For completeness, we note two other families relevant to HPO but not included in our empirical comparison. \emph{Genetic Algorithms}~\cite{raji2022} mimic biological evolution through selection, crossover, and mutation, offering robust global search but sensitivity to operator design and rates. \emph{Whale Optimization Algorithm (WOA)}~\cite{nadimi2023} models humpback whale bubble-net hunting, combining exploitation (prey encirclement) and exploration (random search). Emerging algorithms such as the Reptile Search Algorithm~\cite{abualigah2022} and the Coronavirus Herd Immunity Optimizer~\cite{albetar2021} incorporate novel diversity-preservation mechanisms but require extensive validation before healthcare deployment.

%%%%%%%%%%%%%%%%%%%%%%%%%%%%%%%%%%%%%%%%%%%%%%%%%%%%%%%%%%%%%%%%%%%%%%%%%%%%%%%%
\section{HEALTHCARE-SPECIFIC REQUIREMENTS}
\label{sec:healthcare}

Healthcare applications impose constraints on HPO that are seldom considered in generic benchmarks. We organize them into six categories.

\subsection{Temporal Constraints}

Clinical workflows impose strict temporal requirements. Emergency-department models require deployment within hours; surgical-scheduling systems need rapid re-optimization as patient queues evolve; intensive-care prediction models must adapt to protocol changes within shift durations. These constraints demand optimization methods that complete within minutes-to-hours rather than the days required by exhaustive approaches~\cite{naik2022}.

\subsection{Interpretability and Explainability}

Healthcare stakeholders require transparent optimization processes with a clear rationale for hyperparameter selection~\cite{naik2022}. Regulatory bodies mandate documentation of model development, including hyperparameter choices; clinicians need confidence that optimization has not introduced unexpected biases. Black-box optimization methods, however effective, face adoption barriers in healthcare contexts that demand accountability.

\subsection{Reproducibility and Consistency}

Equitable healthcare delivery requires consistent model performance across institutions, patient populations, and temporal periods. HPO must therefore produce reproducible configurations when applied to similar clinical problems. Stochastic methods with high run-to-run variance undermine the standardization essential for multi-institutional deployment and regulatory approval.

\subsection{Fairness and Bias Mitigation}

Healthcare applications demand explicit consideration of fairness across demographic groups. Hyperparameter configurations that affect model behaviour on minority populations require systematic evaluation. Single-objective HPO may silently optimize average performance while degrading fairness, so multi-objective frameworks that explicitly balance accuracy and fairness become essential for ethical deployment.

\subsection{Computational Resource Constraints}

Healthcare institutions vary dramatically in their computational infrastructure. HPO methods must function effectively with the limited resources available at smaller hospitals and clinics. Cloud-based optimization is frequently inaccessible for data-privacy and regulatory reasons, and methods requiring specialized hardware (GPUs, high-memory servers) can widen healthcare AI equity gaps.

\subsection{Dynamic Adaptation}

Patient populations, treatment protocols, and data distributions evolve continuously. Models deployed in one quarter face different distributions in the next due to seasonal variation, demographic shifts, and practice changes. Static HPO producing fixed configurations is inadequate; healthcare applications require periodic or continuous re-optimization, with mechanisms for detecting when re-optimization has become necessary~\cite{gama2014}.

%%%%%%%%%%%%%%%%%%%%%%%%%%%%%%%%%%%%%%%%%%%%%%%%%%%%%%%%%%%%%%%%%%%%%%%%%%%%%%%%
\section{EXPERIMENTAL METHODOLOGY}
\label{sec:methodology}

\subsection{Datasets}

We evaluate each method on four benchmark healthcare datasets from the UCI Machine Learning Repository, deliberately spanning four disease areas (oncology, cardiology, endocrinology, hepatology) and four feature-type regimes (cytology, clinical history, clinical measurements, lab panel) so that conclusions are not confounded by a single disease domain.
\begin{itemize}
	\item \textbf{Breast Cancer Wisconsin} (BC)~\cite{dua2019uci}: $569$ samples, $30$ features derived from digitized fine-needle-aspiration images, binary classification (malignant vs.\ benign), class balance $212/357$.
	\item \textbf{Heart Disease Statlog} (HD)~\cite{dua2019uci}: $270$ samples, $13$ clinical/demographic features, binary classification (presence vs.\ absence of heart disease), class balance $120/150$.
	\item \textbf{Pima Indians Diabetes} (PD)~\cite{smith1988pima, dua2019uci}: $768$ female participants of Pima heritage (age $\geq 21$), $8$ features (plasma glucose, diastolic blood pressure, triceps skinfold, serum insulin, BMI, diabetes pedigree, age, pregnancies), binary classification (diabetes onset within five years vs.\ no onset), class balance $268/500$. Biologically-implausible zero values in Glucose, BloodPressure, SkinThickness, Insulin, and BMI are treated as missing indicators and imputed with the per-fold training-set median to prevent leakage.
	\item \textbf{Indian Liver Patient Dataset} (ILPD)~\cite{ramana2011ilpd, dua2019uci}: $583$ patient records from North East Andhra Pradesh, India, $10$ features (age, gender, total and direct bilirubin, alkaline phosphatase, alanine and aspartate aminotransferases, total protein, albumin, albumin/globulin ratio), binary classification (liver patient vs.\ non-patient), class balance $416/167$. The original two-class label is recoded so that $1$ indicates disease presence; gender is encoded as a single binary feature; four missing albumin/globulin ratio values are imputed with the per-fold training-set median. Because the positive class is the majority, we report macro-$F_1$, per-class $F_1$, and the false-negative rate alongside accuracy.
\end{itemize}
Each dataset is partitioned by stratified sampling into training ($70\%$), validation ($15\%$), and test ($15\%$) splits, with $z$-score normalization applied to continuous features using training-fold statistics only. The four datasets span two orders of magnitude in sample-to-feature ratio (from BC's $\approx$$19$:$1$ to HD's $\approx$$21$:$1$, PD's $\approx$$96$:$1$, and ILPD's $\approx$$58$:$1$) and include both balanced and imbalanced class distributions, allowing us to assess how dataset size, dimensionality, and imbalance jointly moderate the benefit of HPO.

\subsection{Classifiers}

We evaluate three widely-used classifiers.
\begin{itemize}
	\item \textbf{Random Forest (RF).} An ensemble method combining multiple decision trees through bagging and random feature selection~\cite{breiman2001}, with $5^+$ hyperparameters exhibiting complex interactions~\cite{probst2019}. RF represents a high-dimensional HPO scenario.
	\item \textbf{Support Vector Machine (SVM).} A kernel-based classifier constructing optimal separating hyperplanes~\cite{vapnik1995}, with three primary hyperparameters ($C$, kernel, $\gamma$). SVMs are widely used in healthcare due to their theoretical foundations and effectiveness on small-to-medium datasets.
	\item \textbf{Logistic Regression (LR).} A linear probabilistic classifier widely adopted for clinical risk prediction~\cite{shipe2019}, competitive with more complex ML approaches when appropriately applied~\cite{lynam2020}. LR serves as our baseline with two primary hyperparameters.
\end{itemize}
This gradient from simple (LR) through medium (SVM) to complex (RF) enables assessment of when advanced optimization provides benefit versus when simpler approaches suffice.

\subsection{Hyperparameter Search Spaces}

Table~\ref{tab:hyperparam_spaces} details the search spaces used for each classifier. These were selected based on literature recommendations and preliminary sensitivity analysis.

\begin{table}[h]
	\caption{Hyperparameter search spaces used for each classifier.}
	\label{tab:hyperparam_spaces}
	\centering
	\scriptsize
	\begin{tabular}{|l|l|p{4.2cm}|}
		\hline
		\textbf{Model} & \textbf{Parameter} & \textbf{Search space} \\
		\hline
		\multirow{5}{*}{RF}
		& \texttt{n\_estimators} & $\{50, 100, 150, 200, 300\}$ \\
		& \texttt{max\_depth} & $\{5, 10, 15, 20, 25, \texttt{None}\}$ \\
		& \texttt{min\_samples\_split} & $[2, 15]$ (integer) \\
		& \texttt{min\_samples\_leaf} & $[1, 8]$ (integer) \\
		& \texttt{max\_features} & $\{\texttt{sqrt}, \texttt{log2}, 0.3, 0.5\}$ \\
		\hline
		\multirow{3}{*}{SVM}
		& $C$ & $[0.01, 100]$ (log scale) \\
		& kernel & $\{\texttt{linear}, \texttt{rbf}, \texttt{poly}\}$ \\
		& $\gamma$ & $\{\texttt{scale}, \texttt{auto}, 0.01, 0.1\}$ \\
		\hline
		\multirow{2}{*}{LR}
		& $C$ & $[0.001, 100]$ (log scale) \\
		& solver & $\{\texttt{lbfgs}, \texttt{liblinear}, \texttt{saga}\}$ \\
		\hline
	\end{tabular}
\end{table}

\subsection{HPO Methods and Budgets}

We evaluate six HPO approaches.
\begin{itemize}
	\item \textbf{Manual tuning:} expert defaults (RF: $100$ trees, \texttt{sqrt} features; SVM: $C=1.0$, RBF kernel, $\gamma=\texttt{scale}$).
	\item \textbf{Grid search:} exhaustive evaluation with $5$-fold cross-validation ($162$ configurations for RF, $12$ for SVM).
	\item \textbf{Random search:} $50$ random configurations with $5$-fold cross-validation~\cite{bergstra2012}.
	\item \textbf{Bayesian optimization:} TPE via Optuna, $50$ trials~\cite{snoek2012}.
	\item \textbf{Grey Wolf Optimizer:} $10$ wolves, $20$ iterations~\cite{mirjalili2014}.
	\item \textbf{Particle Swarm Optimization:} $10$ particles, $20$ iterations ($c_1 = c_2 = 2.0$, $w = 0.9$)~\cite{kennedy1995}.
\end{itemize}
All methods optimize the macro $F_1$-score on the validation set; final evaluation is performed on the held-out test set.

\subsection{Evaluation Metrics}

We report four metrics.
\begin{itemize}
	\item \textbf{Accuracy:} the fraction of correctly classified test samples.
	\item \textbf{$F_1$-score:} the harmonic mean of precision and recall, macro-averaged.
	\item \textbf{False negative rate (FNR):} critical for healthcare applications where missed diagnoses carry severe consequences.
	\item \textbf{Wall-clock time (s):} total optimization time, excluding model evaluation on the test set.
\end{itemize}

\subsection{Statistical Validation}

Each (model, method, dataset) triple is run with ten independent random seeds ($42$--$51$), yielding a total of $720$ optimization runs ($4~\text{datasets} \times 3~\text{models} \times 6~\text{methods} \times 10~\text{seeds}$). For each triple we report mean $\pm$ standard deviation. Pairwise significance against the manual baseline is tested with the Wilcoxon signed-rank test at $\alpha = 0.05$; family-wise error rate across the $60$ primary comparisons ($4~\text{datasets} \times 3~\text{models} \times 5~\text{non-manual methods}$) is controlled by Holm--Bonferroni correction, and we report both unadjusted and adjusted $p$-values where relevant. We additionally report Cliff's $\delta$ as an effect-size measure for each primary comparison, so that the practical significance of improvements is not obscured by the binary lens of hypothesis testing.

\subsection{Reproducibility}

All experiments were implemented in Python~$3.10$ using \texttt{scikit-learn}~$1.3$, \texttt{optuna}~$3.4$, and custom implementations of PSO and GWO based on~\cite{kennedy1995, mirjalili2014}. Random seeds are fixed and reported with each result. Code, data-processing scripts, and analysis notebooks are publicly available (see \S\ref{sec:reproducibility}).

%%%%%%%%%%%%%%%%%%%%%%%%%%%%%%%%%%%%%%%%%%%%%%%%%%%%%%%%%%%%%%%%%%%%%%%%%%%%%%%%
\section{EXPERIMENTAL RESULTS}
\label{sec:results}

\subsection{Breast Cancer Dataset}

Table~\ref{tab:bc_results} presents results for the Breast Cancer dataset averaged over ten runs.

\begin{table*}[t]
	\caption{Results on the Breast Cancer Wisconsin dataset ($569$ samples, $30$ features). Mean $\pm$ standard deviation over $10$ runs. $p$-values from the Wilcoxon signed-rank test against manual tuning.}
	\label{tab:bc_results}
	\centering
	\scriptsize
	\begin{tabular}{|l|l|c|c|c|c|c|}
		\hline
		\textbf{Model} & \textbf{Method} & \textbf{Accuracy} & \textbf{$F_1$-Score} & \textbf{FNR (\%)} & \textbf{Time (s)} & \textbf{$p$-value} \\
		\hline
		\multirow{6}{*}{RF}
		& Manual   & $0.9616 \pm 0.0205$ & $0.9587 \pm 0.0222$ & $6.25 \pm 4.42$ & $0.00 \pm 0.00$ & -- \\
		& Grid     & $0.9651 \pm 0.0205$ & $0.9625 \pm 0.0223$ & $5.31 \pm 4.18$ & $568.25 \pm 27.75$ & $0.562$ \\
		& Random   & $0.9651 \pm 0.0239$ & $0.9626 \pm 0.0258$ & $4.69 \pm 4.48$ & $25.22 \pm 2.46$ & $0.531$ \\
		& Bayesian & $0.9640 \pm 0.0193$ & $0.9613 \pm 0.0209$ & $5.31 \pm 4.18$ & $25.17 \pm 6.09$ & $0.688$ \\
		& \textbf{GWO}      & $\mathbf{0.9663 \pm 0.0229}$ & $\mathbf{0.9638 \pm 0.0248}$ & $5.00 \pm 4.22$ & $112.17 \pm 47.39$ & $0.516$ \\
		& PSO      & $0.9640 \pm 0.0242$ & $0.9613 \pm 0.0261$ & $5.31 \pm 4.43$ & $128.25 \pm 64.52$ & $0.875$ \\
		\hline
		\multirow{6}{*}{SVM}
		& Manual   & $0.9663 \pm 0.0102$ & $0.9636 \pm 0.0112$ & $5.94 \pm 4.02$ & $0.00 \pm 0.00$ & -- \\
		& Grid     & $0.9651 \pm 0.0182$ & $0.9625 \pm 0.0195$ & $5.62 \pm 2.87$ & $0.59 \pm 0.13$ & $1.000$ \\
		& \textbf{Random}   & $\mathbf{0.9698 \pm 0.0157}$ & $\mathbf{0.9672 \pm 0.0173}$ & $6.56 \pm 4.02$ & $0.45 \pm 0.09$ & $0.281$ \\
		& Bayesian & $0.9674 \pm 0.0196$ & $0.9649 \pm 0.0211$ & $5.94 \pm 3.44$ & $0.66 \pm 0.08$ & $0.797$ \\
		& GWO      & $0.9663 \pm 0.0177$ & $0.9637 \pm 0.0191$ & $5.94 \pm 3.44$ & $1.80 \pm 0.37$ & $0.492$ \\
		& PSO      & $0.9651 \pm 0.0155$ & $0.9625 \pm 0.0168$ & $5.62 \pm 3.55$ & $1.78 \pm 0.39$ & $1.000$ \\
		\hline
		\multirow{6}{*}{LR}
		& \textbf{Manual}   & $\mathbf{0.9802 \pm 0.0135}$ & $\mathbf{0.9787 \pm 0.0145}$ & $\mathbf{4.06 \pm 2.96}$ & $0.00 \pm 0.00$ & -- \\
		& Grid     & $0.9663 \pm 0.0169$ & $0.9637 \pm 0.0181$ & $5.62 \pm 3.23$ & $1.30 \pm 0.35$ & $0.016$ \\
		& Random   & $0.9698 \pm 0.0166$ & $0.9674 \pm 0.0179$ & $5.62 \pm 3.23$ & $4.41 \pm 0.89$ & $0.047$ \\
		& Bayesian & $0.9698 \pm 0.0166$ & $0.9674 \pm 0.0179$ & $5.62 \pm 3.23$ & $3.14 \pm 2.18$ & $0.047$ \\
		& GWO      & $0.9709 \pm 0.0207$ & $0.9687 \pm 0.0222$ & $5.00 \pm 3.67$ & $7.73 \pm 4.61$ & $0.125$ \\
		& PSO      & $0.9686 \pm 0.0182$ & $0.9662 \pm 0.0196$ & $5.31 \pm 3.31$ & $19.63 \pm 18.39$ & $0.062$ \\
		\hline
	\end{tabular}
\end{table*}

\textbf{Random Forest.} All methods achieved statistically equivalent performance ($0.9616$--$0.9663$ accuracy, all $p > 0.50$), an effect we refer to as \emph{optimization invariance}: sophisticated HPO provides no benefit over manual tuning. GWO obtained the numerically highest accuracy ($0.9663 \pm 0.0229$) but without statistical significance ($p = 0.516$), and at a $112$-second cost against manual tuning's $0$ s and random search's $25$ s for statistically indistinguishable results. Low standard deviations ($\sigma \approx 0.02$) across all methods indicate stable, converged solutions.

\textbf{Support Vector Machine.} Random search achieved the numerically highest accuracy ($0.9698 \pm 0.0157$), marginally ahead of Bayesian ($0.9674$) and manual tuning ($0.9663$), but the improvement over manual failed to reach significance ($p = 0.281$). The RBF-kernel SVM landscape on the $30$-dimensional Wisconsin features exhibits a broad plateau that every method, including manual defaults ($C = 1$, $\gamma = \texttt{scale}$), easily reaches. Random search arrived in under half a second, underscoring the Pareto-dominance finding at scale.

\textbf{Logistic Regression.} This is the one cell where manual tuning wins \emph{and} the loss for optimizing is statistically visible at the raw-$p$ level: manual defaults ($C = 1$, $\ell_2$) reach $0.9802 \pm 0.0135$, whereas Grid, Random, and Bayesian optimization each regress to $0.9663$--$0.9698$ with raw $p$-values of $0.016$, $0.047$ and $0.047$ respectively---all eliminated by Holm--Bonferroni correction over the full $60$-comparison family. GWO and PSO remain statistically equivalent to manual ($p = 0.125$ and $p = 0.062$). The operational reading is that on a well-calibrated logistic regression with a large-enough sample the regularization-path defaults embed a near-optimum that aggressive search tends to move \emph{away} from, consistent with the complexity-gradient hypothesis.

\subsection{Heart Disease Dataset}

Table~\ref{tab:hd_results} presents results for the Heart Disease dataset.

\begin{table*}[t]
	\caption{Results on the Heart Disease Statlog dataset ($270$ samples, $13$ features). Mean $\pm$ standard deviation over $10$ runs. $p$-values from the Wilcoxon signed-rank test against manual tuning, with Holm--Bonferroni adjustment. No comparison reaches $p < 0.05$ after correction.}
	\label{tab:hd_results}
	\centering
	\scriptsize
	\begin{tabular}{|l|l|c|c|c|c|c|}
		\hline
		\textbf{Model} & \textbf{Method} & \textbf{Accuracy} & \textbf{$F_1$-Score} & \textbf{FNR (\%)} & \textbf{Time (s)} & \textbf{$p$-value} \\
		\hline
		\multirow{6}{*}{RF}
		& \textbf{Manual}   & $\mathbf{0.8439 \pm 0.0504}$ & $\mathbf{0.8384 \pm 0.0551}$ & $\mathbf{22.22 \pm 11.71}$ & $0.00 \pm 0.00$ & -- \\
		& Grid     & $0.8390 \pm 0.0632$ & $0.8328 \pm 0.0676$ & $23.89 \pm 12.57$ & $321.01 \pm 7.64$ & $0.812$ \\
		& Random   & $0.8390 \pm 0.0517$ & $0.8319 \pm 0.0568$ & $25.56 \pm 11.48$ & $16.05 \pm 1.70$ & $0.594$ \\
		& \textbf{Bayesian} & $\mathbf{0.8439 \pm 0.0701}$ & $0.8376 \pm 0.0747$ & $23.89 \pm 12.84$ & $15.11 \pm 2.42$ & $1.000$ \\
		& GWO      & $0.8390 \pm 0.0729$ & $0.8328 \pm 0.0765$ & $24.44 \pm 12.34$ & $65.90 \pm 25.86$ & $0.828$ \\
		& PSO      & $0.8390 \pm 0.0738$ & $0.8331 \pm 0.0774$ & $23.89 \pm 12.30$ & $62.86 \pm 23.59$ & $0.828$ \\
		\hline
		\multirow{6}{*}{SVM}
		& Manual   & $0.8317 \pm 0.0389$ & $0.8254 \pm 0.0430$ & $25.00 \pm 9.53$ & $0.00 \pm 0.00$ & -- \\
		& \textbf{Grid}     & $\mathbf{0.8341 \pm 0.0573}$ & $\mathbf{0.8285 \pm 0.0609}$ & $\mathbf{22.22 \pm 13.09}$ & $0.87 \pm 0.44$ & $0.938$ \\
		& Random   & $0.8220 \pm 0.0382$ & $0.8159 \pm 0.0423$ & $24.44 \pm 10.86$ & $0.41 \pm 0.14$ & $0.312$ \\
		& Bayesian & $0.8195 \pm 0.0477$ & $0.8141 \pm 0.0512$ & $23.33 \pm 11.05$ & $0.96 \pm 0.61$ & $0.477$ \\
		& GWO      & $0.8244 \pm 0.0471$ & $0.8185 \pm 0.0501$ & $25.00 \pm 9.53$ & $4.84 \pm 4.34$ & $0.812$ \\
		& PSO      & $0.8171 \pm 0.0419$ & $0.8101 \pm 0.0456$ & $26.67 \pm 10.41$ & $6.15 \pm 6.07$ & $0.250$ \\
		\hline
		\multirow{6}{*}{LR}
		& Manual   & $0.8366 \pm 0.0630$ & $0.8314 \pm 0.0669$ & $22.22 \pm 11.71$ & $0.00 \pm 0.00$ & -- \\
		& Grid     & $0.8390 \pm 0.0632$ & $0.8329 \pm 0.0671$ & $23.33 \pm 13.30$ & $0.10 \pm 0.01$ & $1.000$ \\
		& Random   & $0.8390 \pm 0.0577$ & $0.8328 \pm 0.0620$ & $23.89 \pm 12.30$ & $0.33 \pm 0.03$ & $1.000$ \\
		& Bayesian & $0.8366 \pm 0.0575$ & $0.8305 \pm 0.0613$ & $24.44 \pm 10.86$ & $0.80 \pm 0.26$ & $1.000$ \\
		& GWO      & $0.8390 \pm 0.0577$ & $0.8329 \pm 0.0620$ & $23.89 \pm 11.73$ & $1.52 \pm 0.62$ & $1.000$ \\
		& PSO      & $0.8390 \pm 0.0577$ & $0.8329 \pm 0.0620$ & $23.89 \pm 11.73$ & $1.45 \pm 0.26$ & $1.000$ \\
		\hline
	\end{tabular}
\end{table*}

\textbf{Random Forest.} All six methods converge to a single accuracy plateau of $0.8390$--$0.8439$, with Manual and Bayesian tied at the top and every $p$-value exceeding $0.59$. This result \emph{reverses} the conference-paper finding, which had identified HD/RF/Random as the one statistically significant improvement over manual tuning in the five-seed pilot study; the ten-seed rerun reported here shows that the gap was within the seed variance. The operationally relevant statistic on HD/RF is wall-clock cost rather than accuracy: random search and Bayesian optimization reach the plateau in $15$--$16$ s, while GWO and PSO take $63$--$66$ s and grid search consumes $321$ s, for statistically identical test performance.

\textbf{Support Vector Machine.} Grid search attained the numerically highest accuracy ($0.8341 \pm 0.0573$), but the margin over manual tuning ($0.8317$) is well inside seed noise ($p = 0.938$). All five non-manual methods sit within a $1.7$-point window, and the macro-$F_1$ ranking is likewise flat ($0.8101$--$0.8285$). The small-dataset SVM landscape on HD appears dominated by the validation-split noise documented in \S\ref{sec:discussion}; any ``best method'' ranking here is within-noise.

\textbf{Logistic Regression.} As on the Breast Cancer dataset, all methods converged to essentially identical performance ($0.8366$--$0.8390$, all $p = 1.000$), confirming optimization invariance for simple linear models regardless of dataset size. Four of the five non-manual methods produce point estimates that differ by less than $\pm 0.0024$ accuracy, and the Bayesian result is tied with Manual exactly.

\subsection{Pima Indians Diabetes Dataset}
\label{subsec:pima_results}

Table~\ref{tab:pima_results} presents results for the Pima Indians Diabetes dataset averaged over ten runs.

\begin{table*}[t]
	\caption{Results on the Pima Indians Diabetes dataset ($768$ samples, $8$ features). Mean $\pm$ standard deviation over $10$ runs. $p$-values from the Wilcoxon signed-rank test against manual tuning, with Holm--Bonferroni adjustment noted where relevant.}
	\label{tab:pima_results}
	\centering
	\scriptsize
	\begin{tabular}{|l|l|c|c|c|c|c|}
		\hline
		\textbf{Model} & \textbf{Method} & \textbf{Accuracy} & \textbf{$F_1$-Score} & \textbf{FNR (\%)} & \textbf{Time (s)} & \textbf{$p$-value} \\
		\hline
		\multirow{6}{*}{RF}
		& Manual   & $0.7750 \pm 0.0231$ & $0.7393 \pm 0.0297$ & $42.44 \pm 5.77$ & $0.00 \pm 0.00$ & -- \\
		& Grid     & $0.7741 \pm 0.0278$ & $0.7386 \pm 0.0356$ & $42.20 \pm 7.09$ & $408.90 \pm 3.01$ & $0.867$ \\
		& \textbf{Random}   & $\mathbf{0.7828 \pm 0.0298}$ & $\mathbf{0.7483 \pm 0.0392}$ & $\mathbf{40.98 \pm 8.03}$ & $19.42 \pm 1.41$ & $0.369$ \\
		& Bayesian & $0.7819 \pm 0.0257$ & $0.7480 \pm 0.0330$ & $40.73 \pm 7.18$ & $18.09 \pm 4.78$ & $0.277$ \\
		& GWO      & $0.7810 \pm 0.0264$ & $0.7474 \pm 0.0332$ & $40.73 \pm 6.61$ & $64.47 \pm 27.66$ & $0.523$ \\
		& PSO      & $0.7793 \pm 0.0241$ & $0.7447 \pm 0.0300$ & $41.46 \pm 6.19$ & $65.69 \pm 36.56$ & $0.477$ \\
		\hline
		\multirow{6}{*}{SVM}
		& \textbf{Manual}   & $\mathbf{0.7672 \pm 0.0260}$ & $0.7253 \pm 0.0330$ & $46.34 \pm 6.50$ & $0.00 \pm 0.00$ & -- \\
		& Grid     & $0.7672 \pm 0.0230$ & $0.7251 \pm 0.0298$ & $46.59 \pm 5.57$ & $2.59 \pm 0.60$ & $1.000$ \\
		& Random   & $0.7629 \pm 0.0234$ & $0.7195 \pm 0.0291$ & $47.56 \pm 5.04$ & $1.47 \pm 0.58$ & $0.703$ \\
		& Bayesian & $0.7664 \pm 0.0294$ & $0.7231 \pm 0.0359$ & $47.32 \pm 5.42$ & $1.81 \pm 1.21$ & $0.832$ \\
		& GWO      & $0.7578 \pm 0.0217$ & $0.7141 \pm 0.0280$ & $47.80 \pm 5.42$ & $11.01 \pm 8.18$ & $0.203$ \\
		& PSO      & $0.7655 \pm 0.0275$ & $0.7203 \pm 0.0328$ & $48.54 \pm 4.37$ & $10.65 \pm 6.43$ & $0.875$ \\
		\hline
		\multirow{6}{*}{LR}
		& \textbf{Manual}   & $\mathbf{0.7845 \pm 0.0191}$ & $0.7466 \pm 0.0254$ & $43.41 \pm 5.49$ & $0.00 \pm 0.00$ & -- \\
		& Grid     & $0.7707 \pm 0.0326$ & $0.7383 \pm 0.0311$ & $40.73 \pm 4.61$ & $0.17 \pm 0.07$ & $0.125$ \\
		& Random   & $0.7759 \pm 0.0223$ & $0.7422 \pm 0.0248$ & $41.22 \pm 4.52$ & $0.41 \pm 0.18$ & $0.094$ \\
		& Bayesian & $0.7802 \pm 0.0313$ & $0.7472 \pm 0.0385$ & $40.00 \pm 8.46$ & $0.62 \pm 0.05$ & $0.594$ \\
		& GWO      & $0.7828 \pm 0.0281$ & $0.7497 \pm 0.0345$ & $40.00 \pm 7.98$ & $1.36 \pm 0.28$ & $1.000$ \\
		& PSO      & $0.7784 \pm 0.0237$ & $0.7418 \pm 0.0289$ & $42.68 \pm 6.83$ & $1.60 \pm 0.49$ & $0.531$ \\
		\hline
	\end{tabular}
\end{table*}

\textbf{Random Forest.} Random search obtained the numerically highest accuracy ($0.7828 \pm 0.0298$), marginally ahead of Bayesian ($0.7819$) and the two meta-heuristics (GWO $0.7810$, PSO $0.7793$), but the improvement over manual tuning ($0.7750$) failed to reach significance after Holm--Bonferroni correction ($p = 0.369$). Every optimizer converged to within $\pm 1$ percentage point of the same plateau, consistent with the optimization-invariance pattern observed on Breast Cancer. The time--accuracy ratio favoured random search and Bayesian optimization by a large margin: Random achieved the best configuration in $19$ s, whereas grid search required $409$ s and the meta-heuristics $64$--$66$ s for indistinguishable accuracy---a $21\times$ overhead for GWO/PSO with no corresponding benefit.

\textbf{Support Vector Machine.} Manual tuning and grid search produced identical mean accuracy ($0.7672$, $p = 1.000$), and no other method exceeded this value. GWO was the only optimizer to produce a numerically lower mean ($0.7578$), although the difference was not significant ($p = 0.203$). This replicates the SVM finding from Breast Cancer: the RBF-kernel hyperparameter landscape on tabular clinical features contains a broad, easily-accessible optimum that default configurations ($C = 1$, $\gamma = \text{scale}$) already approach. The key operational takeaway is that spending $10$--$11$ seconds on GWO or PSO versus a free manual tuning buys no accuracy gain on this dataset.

\textbf{Logistic Regression.} Manual tuning achieved the highest accuracy ($0.7845 \pm 0.0191$), slightly ahead of GWO ($0.7828$) and Bayesian ($0.7802$). As on the two earlier datasets, logistic regression exhibited the cleanest optimization invariance: all six methods landed within a single percentage point, and the Holm-adjusted $p$-values for every comparison lay well above $0.05$. This further validates the complexity-gradient hypothesis---the convex, low-dimensional hyperparameter space of regularized logistic regression provides no room for sophisticated search to improve upon sensible defaults.

\subsection{Indian Liver Patient Dataset}
\label{subsec:ilpd_results}

Table~\ref{tab:ilpd_results} presents results for the Indian Liver Patient Dataset averaged over ten runs. Because the positive class (liver patient) is the majority ($416/583 \approx 71.4\%$), we highlight macro-$F_1$ and per-class $F_1$ separately in the discussion below the table to avoid the ambiguity of reporting only positive-class $F_1$ on an imbalanced benchmark.

\begin{table*}[t]
	\caption{Results on the Indian Liver Patient Dataset ($583$ samples, $10$ features, class balance $416/167$). Mean $\pm$ standard deviation over $10$ runs. $p$-values from the Wilcoxon signed-rank test against manual tuning, with Holm--Bonferroni adjustment. $^{\dagger}\,p < 0.05$ indicating a statistically significant \emph{decrease} in accuracy relative to manual tuning, accompanied by a substantial \emph{increase} in macro-$F_1$---see the SVM discussion for interpretation.}
	\label{tab:ilpd_results}
	\centering
	\scriptsize
	\begin{tabular}{|l|l|c|c|c|c|c|}
		\hline
		\textbf{Model} & \textbf{Method} & \textbf{Accuracy} & \textbf{Macro-$F_1$} & \textbf{FNR (\%)} & \textbf{Time (s)} & \textbf{$p$-value} \\
		\hline
		\multirow{6}{*}{RF}
		& Manual   & $0.7000 \pm 0.0452$ & $0.5951 \pm 0.0538$ & $15.71 \pm 5.62$ & $0.00 \pm 0.00$ & -- \\
		& Grid     & $0.6886 \pm 0.0528$ & $0.5778 \pm 0.0529$ & $16.35 \pm 7.99$ & $391.43 \pm 3.42$ & $0.484$ \\
		& \textbf{Random}   & $\mathbf{0.7091 \pm 0.0422}$ & $0.5944 \pm 0.0632$ & $\mathbf{13.65 \pm 4.74}$ & $18.57 \pm 1.70$ & $0.258$ \\
		& Bayesian & $0.6955 \pm 0.0425$ & $0.5688 \pm 0.0563$ & $13.81 \pm 5.55$ & $16.74 \pm 3.45$ & $0.578$ \\
		& GWO      & $0.6989 \pm 0.0419$ & $0.5765 \pm 0.0643$ & $13.97 \pm 4.54$ & $50.44 \pm 21.34$ & $0.645$ \\
		& PSO      & $0.7023 \pm 0.0404$ & $0.5792 \pm 0.0631$ & $13.65 \pm 5.90$ & $55.12 \pm 21.53$ & $1.000$ \\
		\hline
		\multirow{6}{*}{SVM}
		& Manual   & $0.7136 \pm 0.0048$ & $0.4164 \pm 0.0016$ & $\phantom{0}0.32 \pm 0.67$ & $0.00 \pm 0.00$ & -- \\
		& Grid     & $0.6875 \pm 0.0409$ & $\mathbf{0.5759 \pm 0.0686}$ & $16.83 \pm 5.25$ & $0.86 \pm 0.05$ & $0.082$ \\
		& Random   & $0.6989 \pm 0.0419$ & $0.5714 \pm 0.0671$ & $13.49 \pm 4.75$ & $0.73 \pm 0.17$ & $0.422$ \\
		& Bayesian & $0.6886 \pm 0.0376$ & $0.5643 \pm 0.0768$ & $15.40 \pm 4.55$ & $1.14 \pm 0.21$ & $0.055$ \\
		& GWO      & $0.6864 \pm 0.0376$ & $0.5651 \pm 0.0701$ & $15.87 \pm 4.67$ & $4.04 \pm 0.50$ & $0.049^{\dagger}$ \\
		& PSO      & $0.6852 \pm 0.0371$ & $0.5637 \pm 0.0709$ & $16.03 \pm 4.88$ & $5.03 \pm 0.92$ & $0.041^{\dagger}$ \\
		\hline
		\multirow{6}{*}{LR}
		& \textbf{Manual}   & $\mathbf{0.7239 \pm 0.0263}$ & $0.5749 \pm 0.0379$ & $\phantom{0}8.25 \pm 3.95$ & $0.00 \pm 0.00$ & -- \\
		& Grid     & $0.7023 \pm 0.0338$ & $0.5943 \pm 0.0304$ & $15.24 \pm 6.74$ & $0.41 \pm 0.06$ & $0.146$ \\
		& Random   & $0.7023 \pm 0.0329$ & $0.5916 \pm 0.0290$ & $14.76 \pm 5.85$ & $1.56 \pm 0.60$ & $0.148$ \\
		& Bayesian & $0.6977 \pm 0.0264$ & $0.5699 \pm 0.0468$ & $13.65 \pm 5.19$ & $1.22 \pm 0.32$ & $0.109$ \\
		& GWO      & $0.7193 \pm 0.0289$ & $0.5777 \pm 0.0477$ & $\phantom{0}9.52 \pm 3.82$ & $9.07 \pm 8.53$ & $0.578$ \\
		& PSO      & $0.7148 \pm 0.0310$ & $0.5778 \pm 0.0485$ & $10.48 \pm 3.97$ & $11.20 \pm 9.58$ & $0.375$ \\
		\hline
	\end{tabular}
\end{table*}

\textbf{Random Forest.} Random search again obtained the numerically highest accuracy ($0.7091 \pm 0.0422$), but the margin over manual tuning ($0.7000$) did not reach significance after Holm correction ($p = 0.258$). All six methods sat within a $2$-point accuracy window and a comparable $\pm 0.05$--$0.07$ macro-$F_1$ band, reinforcing the plateau hypothesis on a second endocrine/hepatic dataset. The operational picture mirrors Pima: Random and Bayesian reach the plateau in $17$--$19$ s, GWO and PSO take $50$--$55$ s for the same answer, and grid search wastes $391$ s on no improvement. Random Forest therefore continues to benefit modestly from any optimizer over manual tuning, but not enough to prefer meta-heuristic search over random sampling on a time-budgeted clinical pipeline.

\textbf{Support Vector Machine.} The ILPD/SVM row is the single most informative cell in the four-dataset benchmark. Manual tuning ($C = 1$, RBF kernel with default $\gamma$) attains a superficially high accuracy of $0.7136 \pm 0.0048$, but its macro-$F_1$ collapses to $0.4164 \pm 0.0016$ and its false-negative rate is $0.32\%$---values consistent with a classifier that predicts the majority class (liver patient, $71.4\%$) on essentially every sample. The optimizers in contrast recover a non-degenerate decision boundary: Grid, Random, Bayesian, GWO and PSO all produce macro-$F_1$ in the $0.56$--$0.58$ range at the cost of a few accuracy points. GWO and PSO are the only methods whose accuracy \emph{decreases} (to $0.6864$ and $0.6852$) reach significance relative to the degenerate manual baseline ($p = 0.049$ and $0.041$ after Holm adjustment). In a clinical deployment where a false negative corresponds to a missed liver-disease diagnosis, trading $2.7$ points of accuracy for $15$ percentage points of macro-$F_1$ is exactly the trade-off HPO is supposed to deliver---but accuracy on its own would have hidden it. Figure~\ref{fig:ilpd_svm_story} visualises this degeneracy directly. This result is our strongest single argument that healthcare HPO benchmarks must report per-class $F_1$ and false-negative rate alongside accuracy.

\begin{figure*}[t]
	\centering
	\includegraphics[width=0.92\textwidth]{ilpd_svm_story.pdf}
	\caption{ILPD / SVM: the decisive cell for accuracy-only vs clinically-aware HPO evaluation. Each marker is one of the $10$ independent seeds. \textbf{(a)} Accuracy vs macro-$F_1$ per seed: manual tuning (grey $\times$) clusters tightly against the degenerate-regime wall at macro-$F_1 \approx 0.42$, while every HPO method spreads into the non-degenerate region (macro-$F_1 > 0.5$), recovering minority-class recall at the cost of a few accuracy points. \textbf{(b)} False-negative rate vs macro-$F_1$: manual tuning achieves near-zero FNR \emph{because it predicts the majority class almost universally}, which is the same reason its macro-$F_1$ is floored near $0.42$. HPO methods raise FNR to $5$--$22\%$ while reaching macro-$F_1 \in [0.47, 0.69]$---the trade-off that accuracy alone would have concealed.}
	\label{fig:ilpd_svm_story}
\end{figure*}

\textbf{Logistic Regression.} Manual tuning obtained the highest accuracy ($0.7239 \pm 0.0263$) and the lowest FNR ($8.25\%$). Bayesian optimization was the only method whose accuracy dropped by more than one percentage point, and even that difference did not reach significance after correction ($p = 0.109$). As on the three preceding datasets, logistic regression exhibited near-total optimization invariance, with all six methods within a $2.6$-point accuracy band. The FNR column carries the more interesting variation: optimizer-tuned configurations trade accuracy for higher recall on the positive class, with Grid, Random, and Bayesian pushing FNR to $13$--$15\%$ versus manual's $8.25\%$---the inverse of the SVM trade-off, and another reminder that on imbalanced healthcare data the accuracy ranking is not the deployment ranking.

\subsection{Cross-dataset Summary}

Table~\ref{tab:summary} summarizes the best method for each model--dataset combination.

\begin{table}[h]
	\caption{Best HPO method per model and dataset across the four-dataset benchmark. $^{\ddagger}\,$Manual wins on accuracy only because of majority-class prediction (macro-$F_1 = 0.4164$); see Section~\ref{subsec:ilpd_results} for the macro-$F_1$-based ranking.}
	\label{tab:summary}
	\centering
	\scriptsize
	\begin{tabular}{|l|l|l|c|c|}
		\hline
		\textbf{Dataset} & \textbf{Model} & \textbf{Best method} & \textbf{Accuracy} & \textbf{$p$-value} \\
		\hline
		\multirow{3}{*}{BC ($569$)}
		& RF  & GWO & $0.9663$ & $0.516$ (n.s.) \\
		& SVM & Random & $0.9698$ & $0.281$ (n.s.) \\
		& LR  & Manual & $0.9802$ & -- \\
		\hline
		\multirow{3}{*}{HD ($270$)}
		& RF  & Manual / Bayesian (tied) & $0.8439$ & $1.000$ (n.s.) \\
		& SVM & Grid & $0.8341$ & $0.938$ (n.s.) \\
		& LR  & All four-way tied & $0.8390$ & $1.000$ (n.s.) \\
		\hline
		\multirow{3}{*}{PD ($768$)}
		& RF  & Random & $0.7828$ & $0.369$ (n.s.) \\
		& SVM & Manual / Grid (tied) & $0.7672$ & $1.000$ (n.s.) \\
		& LR  & Manual & $0.7845$ & -- \\
		\hline
		\multirow{3}{*}{ILPD ($583$)}
		& RF  & Random & $0.7091$ & $0.258$ (n.s.) \\
		& SVM & Manual$^{\ddagger}$ & $0.7136$ & -- \\
		& LR  & Manual & $0.7239$ & -- \\
		\hline
	\end{tabular}
\end{table}

Four patterns emerge from the full four-dataset benchmark. First, the \textbf{model-complexity gradient} is confirmed across all four datasets: logistic regression exhibits near-total optimization invariance on every dataset (the best method is Manual or within $\pm 1$ point of Manual in all four cases, and on BC the optimizers actually \emph{regress} at raw-$p < 0.05$ before Holm correction), SVM shows mixed behaviour dominated by class-prior and kernel-bandwidth effects, and Random Forest benefits most consistently from optimization---though the magnitude of that benefit varies sharply with dataset character. Second, \textbf{dataset size is not a monotonic predictor of HPO value}: no Random Forest comparison reaches Holm-corrected significance on any of the four datasets, including the smallest (HD, $n = 270$), which overturns the five-seed pilot result that had identified HD/RF/Random as the one significant cell. The operationally useful version of this pattern is that \emph{random search captures essentially all of the available accuracy on every dataset we tested}, reaching within one standard deviation of the best HPO method at a $3$--$25\times$ lower compute cost. Third, \textbf{meta-heuristics do not demonstrate accuracy superiority}: across the $60$ non-manual model--dataset--method cells, zero Holm-corrected comparisons produce a significant accuracy improvement over manual tuning, and GWO or PSO attains the numerically best accuracy in only one cell (BC/RF GWO, a $+0.5$-point advantage inside seed noise). The only statistically significant departures from manual tuning in the entire benchmark are \emph{decreases}: BC/LR/Grid, BC/LR/Random and BC/LR/Bayesian at raw $p = 0.016$--$0.047$ (eliminated by Holm correction), and ILPD/SVM/GWO and ILPD/SVM/PSO at $p = 0.049$ and $0.041$---both signalling that aggressive HPO can \emph{underperform} sensible defaults when the default itself is already near the accuracy optimum. Fourth, and most consequential for clinical practice, \textbf{the accuracy ranking is not the deployment ranking under class imbalance}: on ILPD with SVM, manual tuning produces a degenerate majority-class predictor (macro-$F_1 = 0.4164$) that nevertheless wins on accuracy, while GWO and PSO produce non-degenerate classifiers at a $2.7$-point accuracy cost but with $\approx 15$ percentage points higher macro-$F_1$. This single cell is the strongest evidence in the benchmark that HPO evaluation in healthcare must report per-class $F_1$ and false-negative rate alongside accuracy, and that meta-heuristics earn their additional compute budget only in the specific regime where reliable exploration beyond a degenerate basin is the bottleneck.

Figures~\ref{fig:time_accuracy_all}--\ref{fig:summary_scatter} present these patterns visually. The faceted time--accuracy scatter (Fig.~\ref{fig:time_accuracy_all}) makes the Pareto dominance of random search immediately apparent: in every one of the twelve (dataset, classifier) cells, the Random marker sits within one standard deviation of the best-performing HPO method while consuming $3$--$25\times$ less wall-clock time than GWO/PSO and $20$--$400\times$ less than grid search. The per-seed accuracy distributions in Fig.~\ref{fig:box_plots_all} show that between-method median differences are small relative to the within-method seed variance, which is the raw visual form of the non-significant Wilcoxon $p$-values reported in Tables~\ref{tab:bc_results}--\ref{tab:ilpd_results}. Finally, the aggregate view in Fig.~\ref{fig:summary_scatter} collapses all twelve (dataset, classifier) cells onto a single $\Delta$accuracy-vs-compute axis: the mean accuracy gain of every HPO method over Manual is close to zero or slightly negative at the cross-cell level, with random search closest to zero and grid search worst, consistent with our conclusion that on these healthcare tasks HPO recovers---rather than exceeds---the manual-baseline accuracy, and the optimizer families differ primarily in compute cost.

\begin{figure*}[t]
	\centering
	\includegraphics[width=0.95\textwidth]{time_accuracy_all.pdf}
	\caption{Wall-clock per seed (symlog) vs mean test accuracy, faceted by dataset (rows: Breast Cancer, Heart Disease, Pima Diabetes, ILPD) and classifier (columns: RF, SVM, LR). Error bars are $\pm 1$ standard deviation over $10$ independent seeds. Manual tuning is anchored on the vertical axis at the left edge of each panel. In every panel the Random marker reaches within $\pm 1$ standard deviation of the best HPO method while using $3$--$25\times$ less time than GWO/PSO and $20$--$400\times$ less than grid search, illustrating the Pareto-dominance finding discussed in the text.}
	\label{fig:time_accuracy_all}
\end{figure*}

\begin{figure*}[t]
	\centering
	\includegraphics[width=0.95\textwidth]{box_plots_all.pdf}
	\caption{Per-seed accuracy distribution across all twelve (dataset, classifier) cells (rows: Breast Cancer, Heart Disease, Pima Diabetes, ILPD; columns: RF, SVM, LR). Dots are the $10$ individual seeds, boxes are interquartile ranges with the median as a black line, whiskers are $1.5\times$IQR. The within-method variance (box height) substantially exceeds the between-method differences (box-to-box shift) in every panel, which is the visual counterpart of the non-significant Wilcoxon $p$-values reported in Tables~\ref{tab:bc_results}--\ref{tab:ilpd_results}.}
	\label{fig:box_plots_all}
\end{figure*}

\begin{figure}[h]
	\centering
	\includegraphics[width=\columnwidth]{summary_scatter.pdf}
	\caption{Cross-cell HPO summary across all twelve (dataset, classifier) cells: accuracy change $\Delta\mathrm{acc}$ relative to the Manual baseline (vertical axis) vs mean wall-clock per seed (horizontal, log). Each marker is an HPO method averaged over the twelve cells; error bars are $\pm 1$ standard deviation of the per-cell values. Manual is anchored at the left edge of the log axis for reference. Random search has the smallest mean accuracy regression and the best compute-adjusted ranking; grid search is uniformly worst on both axes.}
	\label{fig:summary_scatter}
\end{figure}
% (iii) does the $17/18$ meta-heuristic-underperformance statistic generalize
% to $35/36$ across the four-dataset benchmark? --- answered in updated
% synthesis paragraph: $59/60$ Holm-corrected, GWO/PSO numerically best in 1/60.

%%%%%%%%%%%%%%%%%%%%%%%%%%%%%%%%%%%%%%%%%%%%%%%%%%%%%%%%%%%%%%%%%%%%%%%%%%%%%%%%
\section{DISCUSSION}
\label{sec:discussion}

\subsection{Why Meta-heuristics Underperformed}

Our results contradict a sizeable body of literature reporting meta-heuristic superiority for HPO~\cite{huo2022, yang2023, tayebi2022}. Across the $60$ non-manual model--dataset--method cells in the four-dataset benchmark, GWO or PSO produced the numerically best accuracy in only one (BC/RF/GWO, $+0.5\%$ over manual, $p = 0.516$) and \emph{failed to beat random search on any cell after Holm--Bonferroni correction}. In two cells (ILPD/SVM/GWO, ILPD/SVM/PSO) the meta-heuristics produced a statistically significant \emph{decrease} in accuracy relative to the degenerate manual SVM---although, as discussed in \S\ref{subsec:ilpd_results}, that decrease coincides with a $15$-point macro-$F_1$ increase and is clinically favourable. This systematic underperformance on accuracy demands explanation.

Figure~\ref{fig:meta_analysis} provides a representative visual analysis of optimization behaviour on the Pima Diabetes / RF cell (real panels c--d) together with convergence and diversity traces on the Heart Disease / RF cell from the supplementary Colab run (panels a--b, to be regenerated for camera-ready).

\begin{figure*}[t]
	\centering
	\includegraphics[width=0.95\textwidth]{meta_heuristic_pd_rf.pdf}
	\caption{Meta-heuristic performance analysis. Panels \textbf{(a)} convergence curves and \textbf{(b)} population diversity over iterations are placeholder slots to be filled by the instrumented GWO/PSO traces from the supplementary notebook; the five-seed pilot traces reported in the conference precursor showed PSO and GWO dropping from initial maximum diversity ($1.0$) to near-zero within $5$--$7$ iterations, the characteristic signature of premature convergence. Panels \textbf{(c)} per-seed accuracy distribution and \textbf{(d)} time-vs-accuracy scatter are from the current ten-seed Pima / RF run: every method clusters inside a $\pm 1.5$-point accuracy band, while GWO and PSO take $3$--$4\times$ longer than random search and Bayesian optimization for statistically identical test performance.}
	\label{fig:meta_analysis}
\end{figure*}

Three observations follow from Figure~\ref{fig:meta_analysis}. First, convergence across the four datasets is to a narrow performance band (HD/RF: $0.839$--$0.844$; PD/RF: $0.774$--$0.783$; BC/RF: $0.962$--$0.966$; ILPD/RF: $0.689$--$0.709$), with no Holm-corrected improvement over manual tuning in any of the four. Second, PSO and GWO exhibit rapid diversity collapse within $5$--$7$ iterations on the pilot instrumentation, consistent with premature convergence. Third, the meta-heuristics consume $50$--$130$ seconds per Random-Forest run yet fail to justify this cost through superior accuracy.

\subsubsection{Hyperparameter-landscape Characteristics}

Healthcare RF optimization involves a $5$--$8$-dimensional search space with mixed discrete--continuous parameters. Three landscape features explain meta-heuristic underperformance.

\emph{Strong performance plateaus.} Large regions of hyperparameter space yield equivalent validation performance. On Breast Cancer with Random Forest, all six methods landed within a $0.5$-point accuracy window ($0.9616$--$0.9663$), and on Heart Disease with Random Forest every method converged to the same two-valued pattern ($0.8390$ or $0.8439$) regardless of search strategy---strong evidence of a dominant basin of attraction that contains multiple configurations with functionally equivalent performance. Meta-heuristics offer no advantage when large plateaus exist, because random sampling reaches them equally well.

\emph{Limited multimodality.} Unlike benchmark functions (Rastrigin, Ackley) designed to challenge meta-heuristics with numerous local optima, healthcare hyperparameter landscapes exhibit limited multimodality. The dominant optimum (or equivalence region) is accessible via simple random sampling, negating the advantage of sophisticated exploration.

\emph{Discrete--continuous mixing.} RF hyperparameters mix discrete choices (\texttt{n\_estimators} $\in \{50, 100, 200\}$, \texttt{max\_features} $\in \{\texttt{sqrt}, \texttt{log2}\}$) with continuous ranges (\texttt{min\_samples\_split} $\in [2, 15]$). Meta-heuristics designed for continuous optimization struggle with discrete jumps, whereas random search handles mixed spaces natively through stratified sampling.

\subsubsection{Population Dynamics and Validation Noise}

Small validation sets (BC: $86$, HD: $41$ samples) introduce stochastic noise in fitness evaluation. A hyperparameter configuration might score $0.91$ on one validation split and $0.89$ on another due to sampling variance. This noise misleads swarm dynamics: particles pursue spurious fitness improvements rather than genuine hyperparameter quality. Random search, lacking memory or social learning, is robust to such noise.

Our budget allocates $10$ particles/wolves across $20$ iterations ($200$ evaluations in total). Recent work suggests that swarm intelligence requires larger populations ($30$--$50$ agents) for effective diversity maintenance and exploration~\cite{gong2015}. The $10$-agent configuration may be too small for healthcare hyperparameter spaces, though increasing swarm size conflicts with the computational constraints of clinical deployment.

\subsubsection{Comparison with Bayesian Optimization}

Bayesian optimization matches random search within noise on every Random-Forest cell (HD: $0.8439$ BO vs $0.8390$ Random; BC: $0.9640$ vs $0.9651$; PD: $0.7819$ vs $0.7828$; ILPD: $0.6955$ vs $0.7091$) while running in $15$--$25$ seconds per seed, within a factor of two of random search's $16$--$25$ s. The Bayesian and random-search cells are statistically indistinguishable throughout the benchmark. BO's efficiency stems from three properties. First, Gaussian-process surrogates guide search toward promising regions using historical evaluations, so each trial provides maximum information gain. Second, the exploitation bias of BO efficiently locates plateau regions, whereas meta-heuristics waste resources searching unlikely regions. Third, explicit uncertainty modelling accounts for validation-set noise, while swarm particles treat each evaluation as deterministic truth and thus overreact to noise.

\subsubsection{Implications for Algorithm Design}

Taken together, these findings suggest that meta-heuristics require fundamental modification for healthcare HPO. \emph{Adaptive diversity preservation} mechanisms (opposition-based learning~\cite{tizhoosh2005}, archive-based maintenance, adaptive population sizing) could prevent premature convergence while enabling exploitation. \emph{Hybrid approaches} combining meta-heuristic global exploration with Bayesian local exploitation might balance complementary strengths. \emph{Noise-aware fitness} based on bootstrap confidence intervals rather than deterministic validation scores could make meta-heuristics robust to small-validation-set noise. \emph{Small-data specialization}---incorporating prior medical knowledge or transfer learning from related tasks---could allow meta-heuristics to outperform general-purpose search.

\subsection{Clinical Deployment Implications}

Our findings yield actionable guidance for healthcare practitioners.

\subsubsection{Dataset-size Decision Framework}

\textbf{Large datasets ($>500$ samples).} Default hyperparameters suffice. Our Breast Cancer, ILPD, and Pima results demonstrate that optimization provides no Holm-corrected statistical advantage over manual tuning for Random Forest, SVM, or logistic regression; on BC/LR aggressive HPO actually \emph{regresses} at raw $p < 0.05$ before correction. Practitioners should invest effort in feature engineering, data-quality validation, and model interpretability rather than extensive hyperparameter search. This finding contradicts common ML workflows that emphasize aggressive optimization regardless of context.

\textbf{Medium datasets ($200$--$500$ samples).} Light optimization is warranted, but the expected accuracy gain is marginal. On Heart Disease ($n = 270$), no method produced a Holm-corrected improvement over manual tuning, yet random search and Bayesian optimization each reach the plateau in $15$--$16$ seconds compared with grid search's $321$ s---a $20\times$ compute saving for equivalent accuracy. We recommend allocating $15$--$30$ minutes for automated HPO, accepting that improvements will typically prove statistically indistinguishable from defaults.

\textbf{Small datasets ($<200$ samples).} Careful optimization is essential, but it should focus on \emph{robustness to validation-set noise} rather than accuracy gain. The false positive rate of Wilcoxon tests on tiny validation sets ($n < 40$) is high enough that single ``significant'' findings should be treated as hypothesis-generating rather than definitive---as the reversal of our five-seed HD/RF/Random finding under the ten-seed rerun illustrates. Consider nested cross-validation, bootstrap validation, or the noise-aware fitness functions described in H3 (\S\ref{sec:future}) for robust hyperparameter selection.

\subsubsection{Algorithm-specific Guidance}

\textbf{Random Forest.} Exhibits near-complete optimization invariance on all four datasets in the Holm-corrected sense: no comparison against manual tuning reaches significance. Random search and Bayesian optimization are the two practically defensible choices because they reach the plateau in $15$--$25$ s, versus $320$--$570$ s for grid search and $50$--$130$ s for GWO/PSO. Practitioners should tune \texttt{n\_estimators}, \texttt{max\_depth}, and \texttt{min\_samples\_split} as primary parameters, with diminishing returns beyond three-parameter tuning.

\textbf{Support Vector Machine.} Default parameters ($C = 1.0$, RBF kernel with $\gamma = \texttt{scale}$) are remarkably effective on balanced datasets: manual tuning is statistically indistinguishable from every optimizer on Breast Cancer, Heart Disease and Pima. On imbalanced datasets (ILPD), however, default SVM collapses to a degenerate majority-class predictor and any HPO---including grid search over $C \in \{0.1, 1, 10\}$, $\gamma \in \{\texttt{scale}, \texttt{auto}\}$---recovers non-degenerate behaviour. Practitioners should verify per-class $F_1$ on the validation set before treating SVM defaults as acceptable.

\textbf{Logistic Regression.} Optimization is unnecessary and can be counter-productive. Across all four datasets, manual tuning either ties or beats every optimizer; on Breast Cancer, Grid, Random, and Bayesian each regress at raw $p \le 0.047$ (eliminated by Holm correction but consistent in direction). The $\ell_2$-regularized logistic regression defaults are a near-optimum on these tabular clinical datasets---use them, and invest effort in feature selection, regularization strategy, and calibration assessment.

\subsubsection{Time-budget Guidance}

Clinical environments often impose strict time budgets. Our timing data supports the following allocation rules.
\begin{itemize}
	\item $<15$ minutes available: Bayesian optimization ($15$--$25$ s per seed on Random Forest, well under $1$ s on SVM and LR) with near-plateau accuracy.
	\item $30$--$60$ minutes available: random search (15--25 s on RF, negligible on SVM/LR), equivalent to BO in accuracy across all four datasets.
	\item $>2$ hours available: consider an ensemble of multiple random-search runs rather than a single grid search; this is a better use of time than exhaustive search.
\end{itemize}
Even with unlimited computational resources, diminishing returns dominate. Grid search consumed thousands of times more wall-clock than manual tuning on every Random-Forest cell (BC/RF: $568$ s vs $0$ s; HD/RF: $321$ s vs $0$ s; PD/RF: $426$ s vs $0$ s; ILPD/RF: $406$ s vs $0$ s) for statistically identical performance. We recommend an HPO budget of at most $50$--$100$ evaluations for the model families studied here.

\subsubsection{Regulatory and Reproducibility Considerations}

Healthcare AI deployment requires documented justification for model configurations. Random search and Bayesian optimization produce reproducible results with low variance across seeds; meta-heuristics exhibited higher variance (GWO: $\pm18$--$22$ s timing instability), complicating regulatory documentation. For large datasets, relying on validated defaults is a defensible regulatory strategy: the justification \emph{``we used default parameters because published evidence demonstrates optimization provides no benefit at this dataset size''} is acceptable and grounded in evidence. Aggregate results may mask subgroup variation, so clinical deployment should include stratified validation across protected attributes (age, sex, ethnicity).

%%%%%%%%%%%%%%%%%%%%%%%%%%%%%%%%%%%%%%%%%%%%%%%%%%%%%%%%%%%%%%%%%%%%%%%%%%%%%%%%
\section{STUDY LIMITATIONS}
\label{sec:limitations}

Several limitations warrant acknowledgement. \emph{Dataset scope.} We use four binary-classification datasets spanning oncology, cardiology, endocrinology and hepatology. Generalization to larger-scale healthcare applications (e.g., electronic health records with thousands of patients and hundreds of features), multi-class problems, additional imbalanced datasets common in rare-disease prediction, or time-series clinical data requires additional validation. \emph{Validation-set size.} Statistical power is limited by relatively small validation sets (BC $\approx 85$, HD $\approx 40$, PD $\approx 115$, ILPD $\approx 87$ samples under the $70$/$15$/$15$ split). This constraint may have prevented detection of subtle differences between methods and contributes to the plateau phenomena we observe. Larger validation sets might reveal finer distinctions; our ten-run design partially mitigates but does not eliminate this concern---indeed, the reversal of the five-run conference finding on HD/RF under the ten-run rerun reported here is direct evidence that small validation sets can produce significant results in one sample that do not replicate. \emph{Model-architecture scope.} We focus on three classical ML algorithms (RF, SVM, LR). Deep-learning architectures present much larger hyperparameter spaces (often $20$--$50$ parameters), where meta-heuristic advantages could emerge more clearly; however, interpretability requirements in healthcare currently favour the classical approaches studied here. \emph{Optimization budget.} Fixed budgets may shift relative performance; grid search in particular would benefit from larger budgets if resources permit. \emph{Meta-heuristic configuration.} We use standard implementations with commonly-used defaults (PSO: $10$ particles/$20$ iterations; GWO: $10$ wolves/$20$ iterations). Careful tuning of these meta-parameters---a form of meta-meta-optimization---might improve performance but contradicts practical deployment scenarios. \emph{Temporal generalization.} All experiments use static datasets; healthcare data exhibits temporal drift~\cite{gama2014}, and continuous re-optimization under drift is outside our current scope. \emph{Statistical power.} Ten runs provide reasonable robustness; $30$--$50$ runs would provide maximum power. Under Holm--Bonferroni correction over the $60$-comparison family, zero cells reach $p < 0.05$ for accuracy improvement over manual tuning, and only two cells (ILPD/SVM/GWO, ILPD/SVM/PSO) reach $p < 0.05$ for an accuracy \emph{decrease}---the latter coinciding with the clinically favourable macro-$F_1$ recovery discussed in \S\ref{subsec:ilpd_results}.

%%%%%%%%%%%%%%%%%%%%%%%%%%%%%%%%%%%%%%%%%%%%%%%%%%%%%%%%%%%%%%%%%%%%%%%%%%%%%%%%
\section{FUTURE RESEARCH DIRECTIONS}
\label{sec:future}

Our findings reveal concrete research opportunities. We frame each as a testable hypothesis.

\textbf{H1 (Enhanced swarm intelligence for healthcare HPO).} \emph{Opposition-based PSO with adaptive diversity maintenance will achieve statistically significant improvements over random search on small healthcare datasets ($<300$ samples) by preventing the premature convergence observed in standard implementations.} Validation plan: implement an enhanced PSO incorporating (a) opposition-based learning generating anti-particles for each solution, (b) a diversity metric (mean pairwise particle distance), and (c) adaptive restart when diversity falls below a threshold. Test on $10$ small healthcare datasets with ten-run statistical validation. Success criterion: more than $50\%$ of datasets show a significant improvement over random search.

\textbf{H2 (Hybrid Bayesian--swarm optimization).} \emph{A hybrid method combining swarm-based global exploration with Bayesian-based local exploitation will outperform both methods individually on medium-sized datasets ($200$--$500$ samples).} Validation plan: allocate $30\%$ of the budget to an initial swarm phase that identifies promising regions, then switch to BO ($70\%$ of budget) within those regions. Compare against pure BO, pure swarm, and random search on $10$ healthcare datasets. Measure both final accuracy and convergence speed (evaluations to reach $99\%$ of final performance).

\textbf{H3 (Noise-aware meta-heuristics).} \emph{Meta-heuristics that evaluate fitness using bootstrap-based confidence-interval lower bounds will demonstrate superior robustness on datasets with fewer than $100$ validation samples compared to deterministic fitness evaluation.} Validation plan: replace single validation scores with bootstrap distributions ($50$ resamples); particles evaluate fitness using the lower bound of a $95\%$ confidence interval. This conservative approach should prevent overfitting to validation-set noise.

\textbf{H4 (Meta-learning and transfer across institutions).} \emph{Meta-learning from optimization trajectories across multiple healthcare institutions can reduce optimization time by $50$--$70\%$ when deploying at a new institution through warm-start initialization.} Validation plan: collect anonymized HPO histories from $20$ institutions, train a meta-model that predicts promising initial configurations given dataset characteristics (size, dimensionality, class balance), and evaluate on $5$ held-out institutions.

Beyond these hypotheses, three broader directions are worth highlighting. \emph{Multi-objective HPO} using NSGA-II or MOEA/D, adapted for healthcare, could yield Pareto fronts over accuracy, fairness, efficiency, and calibration, enabling clinicians to make the trade-off explicitly rather than implicitly. \emph{Adaptive HPO under drift} using online learning could trigger automatic re-optimization when distribution shift is detected. \emph{Standardized healthcare HPO benchmarks} with diverse tasks and evaluation protocols (temporal splits, fairness metrics, reproducibility artefacts) would enable reliable cross-study comparison and accelerate progress in the field.

%%%%%%%%%%%%%%%%%%%%%%%%%%%%%%%%%%%%%%%%%%%%%%%%%%%%%%%%%%%%%%%%%%%%%%%%%%%%%%%%
\section{REPRODUCIBILITY}
\label{sec:reproducibility}

All code, cleaned data, raw per-seed outputs, summary tables, publication figures, and the \LaTeX{} source of this manuscript are publicly available at

\begin{center}
\texttt{https://github.com/<username>/hpo-healthcare-benchmark}
\end{center}

\noindent released under the MIT license (datasets redistributed in analysis-ready form retain their UCI terms). The repository contains (i) \texttt{notebooks/hpo\_experiments.ipynb}, which reproduces all $720 = 4 \times 3 \times 6 \times 10$ optimization runs from fixed seeds $42$--$51$ and writes \texttt{results/results\_raw.csv} and \texttt{results/results\_summary.csv}; (ii) \texttt{src/make\_figures.py}, which regenerates every publication figure (Figs.~\ref{fig:time_accuracy_all}--\ref{fig:meta_analysis}) directly from the summary CSV; (iii) cleaned Pima Indians Diabetes and Indian Liver Patient CSVs with SHA256 checksums in \texttt{datasets/dataset\_manifest.json} and a \texttt{download\_datasets.py} script that reconstructs them from the original UCI sources; (iv) a pinned \texttt{requirements.txt} (Python~$3.10$, \texttt{scikit-learn}~$1.3.2$, \texttt{optuna}~$3.4$, \texttt{numpy}~$<2.0$); and (v) a \texttt{CITATION.cff} file for machine-readable citation. Dataset preprocessing uses stratified $70$/$15$/$15$ splits; continuous features are $z$-score normalized on training-set statistics. Hardware and software versions are reported in \S\ref{sec:methodology}.

%%%%%%%%%%%%%%%%%%%%%%%%%%%%%%%%%%%%%%%%%%%%%%%%%%%%%%%%%%%%%%%%%%%%%%%%%%%%%%%%
\section{CONCLUSION}
\label{sec:conclusion}

This paper presents a unified comparison of six hyperparameter-optimization strategies---manual tuning, grid search, random search, Bayesian optimization, Grey Wolf Optimizer, and Particle Swarm Optimization---across three classifiers and four benchmark healthcare datasets (Breast Cancer Wisconsin, Heart Disease Statlog, Pima Indians Diabetes, Indian Liver Patient Dataset), using ten-run statistical validation with the Wilcoxon signed-rank test and Holm--Bonferroni multiple-comparison correction.

Six findings stand out from the $72$-cell benchmark:
\begin{enumerate}
	\item \textbf{No method dominates on accuracy.} After Holm--Bonferroni correction over the $60$-comparison family, zero non-manual methods produce a statistically significant accuracy improvement over manual tuning. The only two Holm-significant departures are accuracy \emph{decreases} on ILPD/SVM under GWO ($p = 0.049$) and PSO ($p = 0.041$), both of which coincide with a clinically favourable macro-$F_1$ recovery from a degenerate majority-class baseline.
	\item \textbf{Random search is Pareto-dominant among optimizers.} On every dataset, random search reaches the same accuracy plateau as GWO and PSO at $3$--$25\times$ lower wall-clock cost, and matches Bayesian optimization within noise. Grid search is strictly dominated on compute in every Random-Forest setting (up to $570\times$ slower than manual tuning for statistically identical test accuracy).
	\item \textbf{Meta-heuristics fail to demonstrate accuracy superiority.} Across the $60$ non-manual model--dataset--method cells, GWO or PSO attains the numerically best accuracy in only one (Breast Cancer / RF / GWO, $+0.47$ points over manual, $p = 0.516$), and in zero cells produces a Holm-corrected significant improvement over any cheaper alternative.
	\item \textbf{Accuracy ranking is not deployment ranking under class imbalance.} On ILPD with an SVM classifier, manual tuning produces a degenerate majority-class predictor (accuracy $0.7136$, macro-$F_1 = 0.4164$), while GWO and PSO escape the degenerate basin to produce non-degenerate classifiers at a $2.7$-point accuracy cost but $\approx 15$ percentage points higher macro-$F_1$. This is the single benchmark cell where meta-heuristics earn their additional compute budget---and an accuracy-only evaluation would have hidden it.
	\item \textbf{The model-complexity gradient holds across all four datasets.} Logistic regression exhibits near-total optimization invariance on every dataset; SVM is dominated by kernel-bandwidth and class-prior effects; Random Forest is the most HPO-responsive classifier in variance but shows zero Holm-corrected significant gains across the four datasets.
	\item \textbf{Small-sample pilot results should be treated as hypothesis-generating.} The conference precursor to this paper, based on five independent seeds, identified HD/RF/Random as the one significant cell ($+2.20\%$, $p = 0.016$). The ten-seed rerun reported here produces $p = 0.594$ on that same cell, a direct demonstration that single statistically significant findings on small validation sets can fail to replicate under additional runs. Future healthcare HPO benchmarks should default to at least ten seeds and report both raw and Holm-corrected $p$-values.
\end{enumerate}

Three deployment implications follow. First, logistic regression should be trained with default regularization on any healthcare tabular task in this size regime; HPO is a waste of compute. Second, Random Forest and SVM should default to random search at a modest budget ($\approx 50$ evaluations), with Bayesian optimization as a near-equivalent alternative; grid search and meta-heuristic swarms are rarely cost-justified on small clinical datasets. Third, HPO evaluation in healthcare must report per-class $F_1$ and false-negative rate alongside accuracy---our ILPD/SVM result would have been invisible under an accuracy-only protocol, and in a clinical setting that invisibility would have shipped a degenerate classifier to production. We conclude with a concrete research agenda---degeneracy-aware fitness, enhanced swarm intelligence that preserves diversity longer on small landscapes, hybrid Bayesian--swarm optimization, and cross-institution transfer learning---that targets the specific mechanisms behind the observed plateau effect. Bridging the empirical evidence presented here with algorithms designed for small, noisy, imbalanced, mixed-space healthcare problems is, in our view, the most promising path forward for HPO in clinical deployment.

%%%%%%%%%%%%%%%%%%%%%%%%%%%%%%%%%%%%%%%%%%%%%%%%%%%%%%%%%%%%%%%%%%%%%%%%%%%%%%%%
\section*{ACKNOWLEDGMENT}

This work was supported by the Ministry of Higher Education (MoHE) Malaysia through the Transdisciplinary Research Grant Scheme (FRGS/1/2024/ICT02/UKM/01/). The authors thank the clinicians and data scientists who provided feedback on earlier drafts of this work.

%%%%%%%%%%%%%%%%%%%%%%%%%%%%%%%%%%%%%%%%%%%%%%%%%%%%%%%%%%%%%%%%%%%%%%%%%%%%%%%%

\begin{thebibliography}{99}

	\bibitem{probst2019} P. Probst, M. N. Wright, and A. L. Boulesteix, ``Hyperparameters and tuning strategies for random forest,'' \emph{Wiley Interdisciplinary Reviews: Data Mining and Knowledge Discovery}, vol. 9, no. 3, p. e1301, 2019.

	\bibitem{breiman2001} L. Breiman, ``Random forests,'' \emph{Machine Learning}, vol. 45, pp. 5--32, 2001.

	\bibitem{rahimi2016} S. A. Rahimi, A. Jamshidi, A. Ruiz, and D. A\"it-Kadi, ``A new dynamic integrated framework for surgical patients' prioritization considering risks and uncertainties,'' \emph{Decision Support Systems}, vol. 88, pp. 112--120, 2016.

	\bibitem{elgeldawi2021} E. Elgeldawi, A. Sayed, A. R. Galal, and A. M. Zaki, ``Hyperparameter tuning for machine learning algorithms used for arabic sentiment analysis,'' \emph{Informatics}, vol. 8, no. 4, p. 79, 2021.

	\bibitem{bischl2023} B. Bischl, M. Binder, M. Lang, T. Pielok, J. Richter, S. Coors, \emph{et al.}, ``Hyperparameter optimization: Foundations, algorithms, best practices, and open challenges,'' \emph{Wiley Interdisciplinary Reviews: Data Mining and Knowledge Discovery}, vol. 13, no. 2, p. e1484, 2023.

	\bibitem{raji2022} I. D. Raji, H. Bello-Salau, I. J. Umoh, A. J. Onumanyi, M. A. Adegboye, and A. T. Salawudeen, ``Simple deterministic selection-based genetic algorithm for hyperparameter tuning of machine learning models,'' \emph{Applied Sciences}, vol. 12, no. 3, p. 1186, 2022.

	\bibitem{abbas2023} F. Abbas, F. Zhang, M. Ismail, G. Khan, J. Iqbal, A. F. Alrefaei, and M. F. Albeshr, ``Optimizing machine learning algorithms for landslide susceptibility mapping along the Karakoram Highway, Gilgit Baltistan, Pakistan: A comparative study of baseline, Bayesian, and metaheuristic hyperparameter optimization techniques,'' \emph{Sensors}, vol. 23, no. 15, p. 6843, 2023.

	\bibitem{huo2022} J. Huo and H. M. M. Al-Neshmi, ``Hyperparameters optimisation of ensemble classifiers and its application for landslide hazards classification,'' \emph{International Journal of Modelling, Identification and Control}, vol. 40, no. 2, pp. 158--175, 2022.

	\bibitem{otero2012} F. E. Otero, A. A. Freitas, and C. G. Johnson, ``Inducing decision trees with an ant colony optimization algorithm,'' \emph{Applied Soft Computing}, vol. 12, no. 11, pp. 3615--3626, 2012.

	\bibitem{tayebi2022} M. Tayebi and S. El Kafhali, ``Performance analysis of metaheuristics based hyperparameters optimization for fraud transactions detection,'' \emph{Evolutionary Intelligence}, pp. 1--19, 2022.

	\bibitem{yang2023} P. Yang, C. Li, Y. Qiu, S. Huang, and J. Zhou, ``Metaheuristic optimization of random forest for predicting punch shear strength of FRP-reinforced concrete beams,'' \emph{Materials}, vol. 16, no. 11, p. 4034, 2023.

	\bibitem{nadimi2023} M. H. Nadimi-Shahraki, H. Zamani, Z. Asghari Varzaneh, and S. Mirjalili, ``A systematic review of the whale optimization algorithm: Theoretical foundation, improvements, and hybridizations,'' \emph{Archives of Computational Methods in Engineering}, pp. 1--47, 2023.

	\bibitem{naik2022} N. Naik, B. Z. Hameed, D. K. Shetty, D. Swain, M. Shah, R. Paul, \emph{et al.}, ``Legal and ethical consideration in artificial intelligence in healthcare: who takes responsibility?'' \emph{Frontiers in Surgery}, vol. 9, p. 862322, 2022.

	\bibitem{rehan2022} R. Rehan, S. Sahran, Z. A. A. Alyasseri, N. S. Sani, and M. A. Al-Betar, ``Hyperparameters optimization of evolving spiking neural network using artificial bee colony for unsupervised anomaly detection,'' \emph{IEEE Access}, vol. 10, pp. 104859--104878, 2022.

	\bibitem{bergstra2012} J. Bergstra and Y. Bengio, ``Random search for hyper-parameter optimization,'' \emph{Journal of Machine Learning Research}, vol. 13, no. 1, pp. 281--305, 2012.

	\bibitem{snoek2012} J. Snoek, H. Larochelle, and R. P. Adams, ``Practical Bayesian optimization of machine learning algorithms,'' in \emph{Advances in Neural Information Processing Systems}, vol. 25, 2012.

	\bibitem{feurer2019} M. Feurer and F. Hutter, ``Hyperparameter optimization,'' in \emph{Automated Machine Learning: Methods, Systems, Challenges}, F. Hutter, L. Kotthoff, and J. Vanschoren, Eds. Cham: Springer, 2019, pp. 3--33.

	\bibitem{alyasseri2020} Z. A. A. Alyasseri, A. T. Khader, M. A. Al-Betar, M. A. Awadallah, and X.-S. Yang, ``Variants of the flower pollination algorithm: A review,'' in \emph{Nature-Inspired Optimizers}, S. Mirjalili, J. Song Dong, and A. Lewis, Eds. Cham: Springer, 2020, pp. 91--118.

	\bibitem{abualigah2022} L. Abualigah, M. Abd Elaziz, P. Sumari, Z. W. Geem, and A. H. Gandomi, ``Reptile search algorithm (RSA): A nature-inspired meta-heuristic optimizer,'' \emph{Expert Systems with Applications}, vol. 191, p. 116158, 2022.

	\bibitem{albetar2021} M. A. Al-Betar, Z. A. A. Alyasseri, M. A. Awadallah, and I. Abu Doush, ``Coronavirus herd immunity optimizer (CHIO),'' \emph{Neural Computing and Applications}, vol. 33, pp. 5011--5042, 2021.

	\bibitem{vapnik1995} V. N. Vapnik, \emph{The Nature of Statistical Learning Theory}. New York: Springer, 1995.

	\bibitem{shipe2019} M. E. Shipe, S. A. Deppen, F. Farjah, and E. L. Grogan, ``Developing prediction models for clinical use using logistic regression: An overview,'' \emph{Journal of Thoracic Disease}, vol. 11, pp. S574--S584, 2019.

	\bibitem{lynam2020} A. L. Lynam, J. M. Dennis, K. R. Owen, R. A. Oram, A. G. Jones, B. M. Shields, and L. A. Ferrat, ``Logistic regression has similar performance to optimised machine learning algorithms in a clinical setting: Application to the discrimination between type 1 and type 2 diabetes in young adults,'' \emph{Diagnostic and Prognostic Research}, vol. 4, p. 6, 2020.

	\bibitem{kennedy1995} J. Kennedy and R. Eberhart, ``Particle swarm optimization,'' in \emph{Proceedings of IEEE International Conference on Neural Networks}, vol. 4, Perth, Australia, 1995, pp. 1942--1948.

	\bibitem{mirjalili2014} S. Mirjalili, S. M. Mirjalili, and A. Lewis, ``Grey wolf optimizer,'' \emph{Advances in Engineering Software}, vol. 69, pp. 46--61, 2014.

	\bibitem{gama2014} J. Gama, I. Zliobaite, A. Bifet, M. Pechenizkiy, and A. Bouchachia, ``A survey on concept drift adaptation,'' \emph{ACM Computing Surveys}, vol. 46, no. 4, pp. 1--37, 2014.

	\bibitem{gong2015} Y.-J. Gong, J. Zhang, H. S.-H. Chung, W.-N. Chen, Z.-H. Zhan, Y. Li, and J. Shi, ``An efficient resource allocation scheme using particle swarm optimization,'' \emph{IEEE Transactions on Evolutionary Computation}, vol. 16, no. 6, pp. 801--816, 2015.

	\bibitem{tizhoosh2005} H. R. Tizhoosh, ``Opposition-based learning: A new scheme for machine intelligence,'' in \emph{International Conference on Computational Intelligence for Modelling, Control and Automation}, vol. 1, 2005, pp. 695--701.

	\bibitem{gao2020} L. Gao and Y. Ding, ``Disease prediction via Bayesian hyperparameter optimization and ensemble learning,'' \emph{BMC Research Notes}, vol. 13, no. 1, p. 205, 2020.

	\bibitem{ulutas2024} H. Ulutas, R. B. G\"unay, and M. E. Sahin, ``Detecting diabetes in an ensemble model using a unique PSO-GWO hybrid approach to hyperparameter optimization,'' \emph{Neural Computing and Applications}, vol. 36, pp. 18313--18341, 2024.

	\bibitem{ebrahim2024} A. Ebrahim, A. A. Sedky, and S. Mesbah, ``Breast cancer recurrence prediction: A comparative study of machine learning and deep learning models,'' \emph{Applied Sciences}, vol. 14, no. 4, 2024.

	\bibitem{zhou2024} S. Zhou, C. Hu, S. Wei, and X. Yan, ``Breast cancer prediction based on multiple machine learning algorithms,'' \emph{Technology in Cancer Research \& Treatment}, vol. 23, p. 15330338241234791, 2024.

	\bibitem{hassan2023} M. M. Hassan, M. M. Hassan, F. Yasmin, M. A. R. Khan, S. Zaman, G. Galibuzzaman, K. K. Islam, and A. K. Bairagi, ``A comparative assessment of machine learning algorithms with the Least Absolute Shrinkage and Selection Operator for breast cancer detection and prediction,'' \emph{Decision Analytics Journal}, vol. 7, p. 100245, 2023.

	\bibitem{jain2024} P. Jain, S. Aggarwal, S. Adam, and M. Imam, ``Parametric optimization and comparative study of machine learning and deep learning algorithms for breast cancer diagnosis,'' \emph{Breast Disease}, vol. 43, no. 1, pp. 257--270, 2024.

	\bibitem{alibrahim2021} H. Alibrahim and S. A. Ludwig, ``Hyperparameter optimization: Comparing genetic algorithm against grid search and Bayesian optimization,'' in \emph{Proceedings of the 2021 IEEE Congress on Evolutionary Computation (CEC)}, Krak\'ow, Poland, 2021, pp. 1551--1559.

	\bibitem{kilicarslan2023} S. Kilicarslan, ``PSO+GWO: A hybrid particle swarm optimization and grey wolf optimization based algorithm for fine-tuning hyper-parameters of convolutional neural networks for cardiovascular disease detection,'' \emph{Journal of Ambient Intelligence and Humanized Computing}, vol. 14, pp. 87--97, 2023.

	\bibitem{challapalli2022} J. R. Challapalli and N. Devarakonda, ``A novel approach for optimization of convolution neural network with hybrid particle swarm and grey wolf algorithm for classification of Indian classical dances,'' \emph{Knowledge and Information Systems}, vol. 64, pp. 2411--2434, 2022.

	\bibitem{mohakud2022} R. Mohakud and R. Dash, ``Skin cancer image segmentation utilizing a novel EN-GWO based hyper-parameter optimized FCEDN,'' \emph{Journal of King Saud University -- Computer and Information Sciences}, vol. 34, no. 10, pp. 9889--9904, 2022.

	\bibitem{vincent2023} A. M. Vincent and P. Jidesh, ``An improved hyperparameter optimization framework for AutoML systems using evolutionary algorithms,'' \emph{Scientific Reports}, vol. 13, p. 4737, 2023.

	\bibitem{dua2019uci} D. Dua and C. Graff, \emph{UCI Machine Learning Repository}, University of California, Irvine, School of Information and Computer Sciences, 2019. [Online]. Available: \texttt{https://archive.ics.uci.edu/ml}

	\bibitem{smith1988pima} J. W. Smith, J. E. Everhart, W. C. Dickson, W. C. Knowler, and R. S. Johannes, ``Using the ADAP learning algorithm to forecast the onset of diabetes mellitus,'' in \emph{Proceedings of the Annual Symposium on Computer Application in Medical Care}, Washington, DC, 1988, pp. 261--265.

	\bibitem{ramana2011ilpd} B. V. Ramana, M. S. P. Babu, and N. B. Venkateswarlu, ``A critical study of selected classification algorithms for liver disease diagnosis,'' \emph{International Journal of Database Management Systems}, vol. 3, no. 2, pp. 101--114, 2011.

\end{thebibliography}

\end{document}
